%% 10장 부동소수점 수 (새 목차, 2026-11-25)
\chapter{Floating-Point Numbers}
\label{ch:floating-point}

\abstract{컴퓨터의 실수는 거의 예외 없이 IEEE~754의 부동소수점 수이다. 이 장은 그 집합의 간격과 비트 배치, 반올림과 표준 모형, ULP, 정확한 뺄셈, 점진적 언더플로, 특수값, 정수 변환, 다른 형식, 컴파일러의 축약을 다룬다. 구현으로는 stardust가 이웃 수를 찾고 ULP를 재는 법, 특수값 규칙, 포화 변환, FMA 정책을 따라간다.}

Python에서 \code{0.1 + 0.2}는 \code{0.30000000000000004}로 찍히고 \code{0.1 + 0.2 == 0.3}은 \code{False}이다. 버그처럼 보이지만 IEEE~754가 정한 그대로의 동작이다. 이 장을 다 읽으면 이 결과를 손으로 예측할 수 있다. $0.1$은 이진법으로 유한하게 적히지 않으므로 리터럴 \code{0.1}은 $0.1$에 가장 가까운 binary64 수, 정확히
\[
\begin{aligned}
\frac{3602879701896397}{2^{55}}&=0.1000000000000000055511151231257827\\
&\phantom{{}=0.1000000000000000}021181583404541015625
\end{aligned}
\]
를 가리킨다. \code{0.2}와 \code{0.3}도 가장 가까운 수로 바뀐다. 두 근삿값의 정확한 합은 다시 가장 가까운 수로 반올림되는데 그 수가 하필 \code{0.3}이 가리키는 수의 바로 다음 수이다. 모든 단계가 ``가장 가까운 수''라는 하나의 규칙을 따랐을 뿐인데 결과가 어긋났다. 수치 계산의 정확도를 따지는 모든 논의는 이 규칙을 정확히 이해하는 데서 시작한다.

\section{The Set of Floating-Point Numbers}
\label{sec:fp-set}

\subsection{Definition}

과학적 표기법은 수를 $6.022\times10^{23}$처럼 유효숫자와 지수로 나누어 적는다. 유효숫자의 자릿수를 고정하면 아주 큰 수와 아주 작은 수를 같은 상대 정밀도로 적을 수 있다. 부동소수점 수는 이 생각을 기수 $2$로 옮긴 것이다.

\begin{definition}[부동소수점 수]
\label{def:fp-set}
정수 $\beta\ge2$(기수), $p\ge1$(정밀도), $e_{\min}<e_{\max}$에 대해
\[
x=\pm\,m\cdot\beta^{e-p+1},\qquad m\in\mathbb Z,\ 0\le m<\beta^p,\ e_{\min}\le e\le e_{\max}
\]
꼴로 쓸 수 있는 실수 $x$들의 집합을 $\mathbb F$라 한다. $m$을 \term{유효숫자}{significand}, $e$를 \term{지수}{exponent}라 한다. $m\ge\beta^{p-1}$로 쓸 수 있는 $x$를 \term{정규 수}{normal number}, $0$이 아니면서 그렇게 쓸 수 없는 $x$를 \term{서브노멀 수}{subnormal number}라 한다.
\end{definition}

유효숫자를 $\beta$진법 $p$자리로 적고 맨 앞자리 뒤에 소수점을 찍으면 $x=\pm d_0.d_1\cdots d_{p-1}\times\beta^e$가 되고 정규 수는 맨 앞자리가 $0$이 아닌 수이다. 서브노멀 수는 지수가 이미 $e_{\min}$이라 앞자리의 $0$을 지수로 옮겨 없앨 수 없는 수이다. 이 책은 거의 언제나 $\beta=2$, $p=53$, $e_{\min}=-1022$, $e_{\max}=1023$인 \textbf{binary64}를 다룬다. C와 C++의 \code{double}, Python의 \code{float}, numpy와 stardust의 \code{float64}가 모두 이 형식이다\cite{ieee754-2019}.

\subsection{The Spacing of Floating-Point Numbers}

$\mathbb F$의 수들은 고르게 퍼져 있지 않다. 지수가 $e$인 정규 수는 $2^e\le|x|<2^{e+1}$에 놓이는데 이 구간을 \term{이진 구간}{binade}이라 부르자. 한 이진 구간 안에서는 유효숫자가 $2^{p-1}$부터 $2^p-1$까지 $1$씩 커지므로 이웃한 두 수의 간격이 일정하게 $2^{e-p+1}$이고 이진 구간을 하나 올라갈 때마다 간격이 두 배가 된다. 그래서 수들은 $0$ 근처에 촘촘하고 멀어질수록 성기다.

간격을 $x$로 나눈 상대 간격은 이진 구간의 아래 끝 $2^e$에서 $2^{1-p}$이고 위 끝 근처에서 그 절반 $2^{-p}$에 가깝다(\cref{fig:spacing}). 상대 정밀도가 한 이진 구간 안에서 두 배까지 출렁이는 이 현상을 \term{흔들림}{wobble}이라 한다. 기수 $\beta$에서는 흔들림이 $\beta$배이므로 상대 정밀도는 기수 $2$에서 가장 고르다. 이진 부동소수점을 쓰는 이유 가운데 하나가 여기에 있다\cite{goldberg1991}.

\begin{figure}[htbp]
\centering
\begin{tikzpicture}
\begin{semilogxaxis}[sdplot,xlabel={$x$},ylabel={$\operatorname{ulp}(x)/x$},xmin=0.25,xmax=16,ymin=0,ymax=2.5e-16,width=0.7\textwidth,height=0.4\textwidth,log basis x=2,scaled y ticks=false,yticklabel style={/pgf/number format/sci,/pgf/number format/precision=1}]
\addplot[black,thick] table[x=x,y=rel]{figures/10/spacing.dat};
\addlegendentry{상대 간격}
\addplot[black!50,dashed,domain=0.25:16] {2^(-52)};
\addlegendentry{$2^{-52}$}
\addplot[black!50,dotted,thick,domain=0.25:16] {2^(-53)};
\addlegendentry{$2^{-53}$}
\end{semilogxaxis}
\end{tikzpicture}
\caption{binary64의 상대 간격. 이진 구간마다 $2^{-52}$에서 $2^{-53}$ 가까이로 줄었다가 다음 구간에서 다시 뛴다(흔들림; \code{figures/10/fp.py}).}
\label{fig:spacing}
\end{figure}

\begin{proposition}
\label{prop:fp-integers}
binary64에서 $|n|\le2^{53}$인 모든 정수 $n$은 정확히 나타난다. $2^{53}+1$은 나타나지 않는다.
\end{proposition}

\begin{proof}
$|n|<2^{53}$이면 $m=|n|$, $e=p-1=52$로 두면 $n=\pm m\cdot2^0$이다. $n=\pm2^{53}$은 $m=2^{52}$, $e=53$으로 쓴다. 한편 $2^{53}\le x<2^{54}$의 간격은 $2^{53-52}=2$이므로 이 구간의 수는 모두 짝수이고 홀수 $2^{53}+1$은 그 가운데 없다.
\end{proof}

$2^{53}=9007199254740992$이다. \code{int64}의 범위는 $-2^{63}$부터 $2^{63}-1$까지이므로 큰 정수를 \code{float64}로 바꾸면 끝자리가 사라질 수 있다(\cref{sec:conversion}).

\subsection{The Bit Layout of binary64}

binary64 수 하나는 $64$비트에 이렇게 담긴다. 맨 위 $1$비트는 부호 $s$, 그다음 $11$비트는 \term{치우친 지수}{biased exponent} $E$, 나머지 $52$비트는 \term{가수 필드}{fraction field} $f$이다. 해석은 $E$에 따라 셋으로 나뉜다.
\begin{enumerate}
\item $1\le E\le2046$이면 정규 수 $(-1)^s(1+f/2^{52})\cdot2^{E-1023}$이다. 정규 수의 맨 앞자리는 이진법에서 언제나 $1$이므로 그 비트는 저장하지 않고 이 \term{숨은 비트}{hidden bit} 덕분에 $52$비트로 $53$자리의 정밀도를 얻는다.
\item $E=0$이면 $(-1)^s(f/2^{52})\cdot2^{-1022}$이다. $f=0$이면 $\pm0$이고 그 밖에는 서브노멀 수이다.
\item $E=2047$이면 $f=0$일 때 $\pm\infty$, $f\ne0$일 때 NaN이다.
\end{enumerate}
\begin{figure}[htbp]
\centering
\begin{tikzpicture}[x=0.17cm,y=0.6cm]
\draw[fill=black!25] (0,0) rectangle (1,1);
\draw[fill=black!10] (1,0) rectangle (12,1);
\draw (12,0) rectangle (64,1);
\foreach \i in {1,...,63} \draw[black!40] (\i,0)--(\i,0.25);
\node[font=\small] at (0.5,1.45) {$s$};
\node[font=\small] at (6.5,1.45) {$E$ ($11$비트)};
\node[font=\small] at (38,1.45) {$f$ ($52$비트)};
\node[font=\scriptsize,anchor=north west,inner xsep=0pt] at (0,-0.05) {63};
\node[font=\scriptsize,anchor=north east,inner xsep=1pt] at (12,-0.05) {52};
\node[font=\scriptsize,anchor=north west,inner xsep=1pt] at (12,-0.05) {51};
\node[font=\scriptsize,anchor=north east,inner xsep=0pt] at (64,-0.05) {0};
\node[font=\small,anchor=north] at (32,-0.8) {$0.1\approx$ \code{0x3FB999999999999A}: $s=0$, $E=1019$, $f=\code{0x999999999999A}$};
\end{tikzpicture}
\caption{binary64의 비트 배치. 맨 위 비트부터 부호, 치우친 지수, 가수 필드이고 아래 숫자는 비트 번호이다. $0.1$의 근삿값은 $(1+f/2^{52})\cdot2^{1019-1023}$이다.}
\label{fig:bit-layout}
\end{figure}

지수에서 $1023$을 빼는 것은 지수를 부호 없는 정수로 적기 위한 장치이다. 그 결과로 다음의 중요한 성질이 생기는데 1장에서 부동소수점 수를 정수 키로 정렬할 때 가정으로 받아들인 성질이다.

\begin{proposition}
\label{prop:fp-bits-monotone}
NaN이 아닌 binary64 수 $x$의 하위 $63$비트를 부호 없는 정수로 읽은 값을 $\mu(x)=2^{52}E+f$라 하자. NaN이 아닌 두 수 $x$, $y$에 대해 $|x|<|y|$와 $\mu(x)<\mu(y)$는 동치이다. 또 $\mu$가 $1$ 커지는 것은 $|x|$가 바로 다음 binary64 수로 가는 것과 같다.
\end{proposition}

\begin{proof}
$\mu$를 $0,1,2,\dots$로 늘려 가며 나타내는 수를 보자. $E=0$인 동안에는 $|x|=f\cdot2^{-1074}$이 $2^{-1074}$씩 커진다. $f=2^{52}-1$에서 $f=0$, $E=1$로 넘어가면 $|x|$는 $(2^{52}-1)2^{-1074}$에서 $2^{-1022}=2^{52}\cdot2^{-1074}$로 역시 $2^{-1074}$만큼 커진다. $E\ge1$인 이진 구간 안에서는 $|x|=(2^{52}+f)\cdot2^{E-1075}$가 일정한 간격 $2^{E-1075}$로 커지고 다음 $E$로 넘어갈 때도 $(2^{53}-1)2^{E-1075}$에서 $2^{52}\cdot2^{E+1-1075}$로 같은 간격만큼 커진다. 마지막 유한 수 다음의 $E=2047$, $f=0$은 $+\infty$이다. 따라서 $\mu\mapsto|x|$는 강한 증가 함수이고 연속한 $\mu$는 이웃한 수에 대응한다.
\end{proof}

이 명제에 따르면 이웃한 수로 가는 연산은 비트 패턴에 $\pm1$을 하는 일로 줄어든다. $0$에서 멀어지는 방향이면 크기가 커지므로 $1$을 더하고 $0$ 쪽으로 가면 $1$을 뺀다(\cref{alg:nextafter}). $0$을 지날 때와 NaN, 같은 두 수만 따로 처리하면 된다.

\begin{algorithm}[htbp]
\begin{algorithmic}[1]
\Require binary64 수 $a$, $b$
\Ensure $a$에서 $b$ 쪽으로 바로 다음 binary64 수
\State \textbf{if} $a$ 또는 $b$가 NaN \textbf{then} \Return $a+b$ \Comment{NaN을 전한다}
\State \textbf{if} $a=b$ \textbf{then} \Return $b$ \Comment{\code{nextafter(0.0, -0.0)}은 $-0$}
\State \textbf{if} $a=0$ \textbf{then} \Return $b$의 부호를 붙인 $2^{-1074}$
\State $n\gets a$의 비트를 부호 있는 $64$비트 정수로 읽은 값
\State \textbf{if} ($a<b$) $=$ ($a>0$) \textbf{then} $n\gets n+1$ \textbf{else} $n\gets n-1$ \Comment{$0$에서 멀어지면 $+1$}
\State \Return $n$을 binary64로 읽은 값 \Comment{가장 큰 유한 수 다음은 $\infty$}
\end{algorithmic}
\caption{비트 패턴으로 걷는 \code{nextafter} (stardust \code{scalar\_ops.hpp}의 \code{Nextafter})}
\label{alg:nextafter}
\end{algorithm}

\cref{alg:nextafter}는 libm의 \code{nextafter}와 같은 결과를 내면서 분기 몇 개와 정수 덧셈 하나로 끝난다. 원소마다 libm을 부르는 numpy와 견주면 $10^6$개 원소에서 한 스레드로 $581\,\mu$s 대 $1450\,\mu$s이다. 기본 스레드 설정($10$개)에서는 $78\,\mu$s이다(복제본, \code{figures/10/review\_sd.py}). 코드 주석이 적은 ``약 $8$배''보다는 차이가 작다. 부호 있는 정수로 읽어도 된다. 음수의 비트 패턴도 부호 비트 아래로는 $\mu(|a|)$이고 $\pm1$이 부호 비트까지 넘어가지 않기 때문이다. $0$은 3행에서 미리 걸렀고 가장 큰 유한 수에 $1$을 더하면 $E=2047$, $f=0$, 곧 $\infty$가 되어 표준의 넘침 규칙과 맞는다. 거꾸로 $+\infty$에서 $0$ 쪽으로 가면 비트에서 $1$을 빼므로 가장 큰 유한 수가 나온다.

$E$와 $f$의 범위를 대입하면 binary64의 주요 상수가 나온다(\cref{tab:binary64}, 정확한 유리수 산술로 확인).

\begin{table}[htbp]
\begin{sdtable}{colspec={lll}, row{1}={font=\small\bfseries}}
\toprule
이름 & 정확한 값 & 십진 근삿값 \\
\midrule
가장 큰 유한 수 & $(2-2^{-52})\cdot2^{1023}$ & $1.7976931348623157\times10^{308}$ \\
가장 작은 양의 정규 수 & $2^{-1022}$ & $2.2250738585072014\times10^{-308}$ \\
가장 작은 양의 서브노멀 수 & $2^{-1074}$ & $4.9406564584124654\times10^{-324}$ \\
기계 엡실론 $\epsilon$ & $2^{-52}$ & $2.220446049250313\times10^{-16}$ \\
단위 반올림 $u$ & $2^{-53}$ & $1.1102230246251565\times10^{-16}$ \\
\bottomrule
\end{sdtable}
\caption{binary64의 주요 상수. 기계 엡실론은 $1$과 그다음 수의 간격이고 단위 반올림은 그 절반이다.}
\label{tab:binary64}
\end{table}

\subsection{Bit-Level Predicates and Exponent Splitting}
\label{sec:bit-ops}

\cref{prop:fp-bits-monotone}과 비트 배치를 알면 부호와 분류와 지수에 관한 연산 대부분이 정수 연산 몇 개로 줄어든다. 부호 비트를 $S=2^{63}$, 지수 필드가 모두 $1$인 패턴을 $\mathtt{0x7FF0\ldots0}$이라 하자. 비트를 왼쪽으로 한 칸 밀면 부호가 사라지고 지수와 가수가 위로 붙으므로 $\pm\infty$는 정확히 $\mathtt{0xFFE0\ldots0}$이 된다. 이 값보다 작으면 유한 수이고 크면 NaN이다(\cref{alg:bit-predicates}). 부호에 관한 연산은 맨 위 비트 하나만 건드린다. 절댓값은 부호 비트를 지우고 부호 바꾸기는 그 비트를 뒤집는다. \code{copysign}은 한쪽의 크기와 다른 쪽의 부호를 이어 붙인다. 이 연산들은 NaN에도 정의되어 있고 비트를 그대로 옮기므로 예외를 일으키지 않는다. 그래서 표준은 이것들을 산술 연산과 따로 묶어 ``조용한 연산''이라 부른다\cite{ieee754-2019}.

\begin{algorithm}[htbp]
\begin{algorithmic}[1]
\Require binary64 수 $x$, $y$의 비트 $b$, $c$ (부호 없는 $64$비트), $S=2^{63}$
\State $\textproc{IsFinite}\gets(b\ll1)<\mathtt{0xFFE0000000000000}$ \Comment{\code{IsFinite}}
\State $\textproc{IsInf}\gets(b\ll1)=\mathtt{0xFFE0000000000000}$ \Comment{\code{IsInf}}
\State $\textproc{IsNaN}\gets(b\ll1)>\mathtt{0xFFE0000000000000}$
\State $\textproc{SignBit}\gets b\gg63$; $\textproc{Abs}\gets b\wedge(S-1)$; $\textproc{Neg}\gets b\oplus S$
\State $\textproc{CopySign}(x,y)\gets(b\wedge(S-1))\vee(c\wedge S)$
\State $\textproc{IsPositiveNormal}\gets(b-\mathtt{0x0010000000000000})<\mathtt{0x7FE0000000000000}$ \Comment{\code{log\_ok}}
\end{algorithmic}
\caption{비트로 하는 분류와 부호 연산 (stardust \code{scalar\_ops.hpp}, \code{vmath.hpp}). $\ll$, $\gg$는 논리 밀기이고 비교는 부호 없는 정수의 비교이다.}
\label{alg:bit-predicates}
\end{algorithm}

마지막 줄은 부호 없는 정수의 감김을 이용한다. 양의 정규 수의 비트는 $\mathtt{0x0010\ldots0}$부터 $\mathtt{0x7FEF\ldots F}$까지이므로 아래 끝을 빼면 $[0,\mathtt{0x7FE0\ldots0})$에 든다. $+0$과 양의 서브노멀 수는 아래 끝보다 작아서 빼면 $2^{64}$ 근처로 감긴다. 음수는 부호 비트 때문에 이미 크고 $+\infty$와 NaN은 위 끝 이상이다. 그래서 비교 한 번이 ``양의 정규 유한 수''를 가려낸다. stardust의 NEON $\log$는 이 판정으로 고른 원소만 벡터 경로로 보내고 나머지는 libm으로 넘긴다. 같은 목적의 다른 관용구도 있다. NaN은 모든 비교에서 거짓이므로 \code{!(fabs(x) <= 708)}처럼 부정한 비교를 쓰면 NaN이 저절로 예외 경로에 떨어진다. stardust의 \code{isinf}는 처음에 \code{std::isinf}를 썼다가 이 비트 판정으로 바꾸었다. 코드 주석에 따르면 \code{std::isinf}는 $10^6$개에 $87$--$94\,\mu$s였다. 지금은 한 스레드로 $65.0\,\mu$s여서 numpy($65.1\,\mu$s)와 같고 기본 스레드 설정에서는 $17.2\,\mu$s이다(복제본).

지수를 떼어 내는 일도 비트로 한다. $x$를 $2^kz$로 나누고 $z$를 어떤 구간 $[c,2c)$에 두고 싶다고 하자($\log$의 인자 축소에서는 $z$가 $1$ 근처에 오도록 $c<1$로 잡는다). $c$의 비트를 빼고 산술 밀기로 $52$비트 밀면 $k$가 나온다. $x$의 비트에서 $k\ll52$를 빼면 $z$가 나온다(\cref{alg:split-exponent}). 증명은 \cref{prop:fp-bits-monotone}에서 바로 나온다. 정규 수의 비트에서 $k\ll52$를 빼면 가수는 그대로 두고 지수에서만 $k$를 빼게 되므로 $2^{-k}$를 곱한 것과 같다. 그리고 $x\ge c$이면 $\mathrm{bits}(x)-\mathrm{bits}(c)$의 지수 자리는 $x/c$가 몇 번째 이진 구간인지를 센다. 반대로 $2^k$를 곱하려면 비트에 $k\ll52$를 더하면 되는데 결과가 정규 수로 남는 동안만 맞다. 그래서 stardust의 $\exp$는 $|x|\le708$일 때만 이 덧셈을 쓴다($e^{\pm708}$과 그 앞뒤의 $2^{k+j/128}$이 정규 수로 남는다).

\begin{algorithm}[htbp]
\begin{algorithmic}[1]
\Require 양의 정규 수 $x$, 상수 $c$ (stardust의 $\log$는 $c=0.685546875$, 비트 \code{0x3FE5F00000000000})
\Ensure $x=2^kz$, $z\in[c,2c)$, 표의 칸 번호 $i$
\State $t\gets\mathrm{bits}(x)-\mathrm{bits}(c)$ \Comment{부호 없는 뺄셈}
\State $k\gets t\gg_{\mathrm{a}}52$ \Comment{산술 밀기: $x<c$이면 음수}
\State $z\gets$ $\mathrm{bits}(x)-(k\ll52)$를 binary64로 읽은 값
\State $i\gets(t\gg45)\wedge127$ \Comment{가수의 위 $7$비트로 $[c,2c)$를 $128$칸으로}
\Function{ScaleByPow2}{$z$, $k$} \Comment{결과가 정규 수일 때만}
  \State \Return $\mathrm{bits}(z)+(k\ll52)$를 binary64로 읽은 값
\EndFunction
\end{algorithmic}
\caption{지수 떼어 내기와 $2$의 거듭제곱 곱하기 (stardust \code{vmath.hpp}의 \code{log\_from}, \code{exp\_parts})}
\label{alg:split-exponent}
\end{algorithm}

\section{Rounding}
\label{sec:rounding}

\subsection{Rounding Functions}

실수 $x$를 $\mathbb F$의 원소로 바꾸는 함수를 반올림이라 한다. IEEE~754의 기본값은 \term{짝수로의 반올림}{round to nearest, ties to even}이다. 가장 가까운 원소로 보내되 두 원소의 한가운데에 있으면 유효숫자가 짝수인 쪽, 곧 마지막 비트가 $0$인 쪽으로 보낸다. 이 밖에 $0$ 쪽, $+\infty$ 쪽, $-\infty$ 쪽으로 자르는 세 가지 \term{방향 반올림}{directed rounding}이 있고 구간 산술처럼 참값을 위아래로 확실히 묶어야 하는 곳에서 쓴다. 별다른 말이 없으면 이 책의 반올림은 짝수로의 반올림이고 이것을 $\fl(x)$로 적는다.

$\mathbb F$의 범위 밖도 정해야 한다. 표준은 지수의 상한이 없는 가상의 집합에서 반올림한 결과의 크기가 가장 큰 유한 수를 넘으면 $\pm\infty$를 결과로 한다. 짝수로의 반올림에서는 $|x|\ge(2-2^{-53})\cdot2^{1023}$, 곧 가장 큰 유한 수와 가상의 다음 수 $2^{1024}$의 한가운데 이상일 때이고 이것을 \term{넘침}{overflow}이라 한다.

짝수 쪽으로 가르는 규칙은 한가운데 값이 생각보다 자주 나오기 때문에 필요하다. 두 수를 더하거나 정수를 바꿀 때 버리는 부분이 정확히 마지막 자리의 절반인 경우는 드물지 않고 그때마다 늘 위로 올리면 오차가 한 방향으로 쌓인다. 짝수로 가르면 한가운데 값의 절반 정도는 위로, 절반 정도는 아래로 가므로 이런 치우침이 사라진다(\cref{fig:bias}). Knuth는 이 규칙에서 $u\gets(u\oplus v)\ominus v$를 거듭해도 첫 단계 이후로는 값이 더 바뀌지 않음을 보였다(\cite{knuth1997taocp2}의 4.2.2절). 한가운데를 늘 위로(0에서 먼 쪽으로) 올리는 규칙에서는 같은 계산이 한 방향으로 흘러갈 수 있다. 정밀도 $4$에서 $u=1/4$, $v=9/64$로 시작하면 $u$가 $9/32$, $5/16$, $11/32$로 커지다가 $3/8$에서 멈춘다. 첫 단계에서 멈추는 짝수 규칙과 달리 여러 번 흘러간다. 정밀도 $4$의 모든 쌍을 검사하면 짝수 규칙에서는 어긋나는 쌍이 없고 위로 올리는 규칙에서는 $1022$쌍이 흘러간다(\code{figures/10/knuth.py}). 정수의 예를 보면 $2^{53}+1$은 $2^{53}$과 $2^{53}+2$의 한가운데이고 유효숫자가 짝수인 쪽은 $2^{53}$이므로 $\fl(2^{53}+1)=2^{53}$이다. 반면 $2^{53}+3$은 $2^{53}+2$와 $2^{53}+4$의 한가운데이고 $2^{53}+4$ 쪽이 짝수이므로 위로 올라간다.

\begin{figure}[htbp]
\centering
\begin{tikzpicture}
\begin{axis}[sdplot,xlabel={반올림한 수의 개수},ylabel={누적 오차},xmin=0,xmax=2000,width=0.62\textwidth,height=0.4\textwidth]
\addplot[black,thick] table[x=k,y=up]{figures/10/bias.dat};
\addlegendentry{한가운데를 위로}
\addplot[black!55,thick] table[x=k,y=even]{figures/10/bias.dat};
\addlegendentry{한가운데를 짝수로}
\end{axis}
\end{tikzpicture}
\caption{$0$에서 $100$까지 $0.5$ 간격의 무작위 수 $2000$개(그 절반쯤은 정수이다)를 정수로 반올림한 오차의 누적(정확한 유리수 계산; \code{figures/10/fp.py}). 위로 올리면 개수에 비례해 $508.5$까지 쌓이고 짝수로 가르면 $11.5$에 머문다.}
\label{fig:bias}
\end{figure}

\subsection{Correctly Rounded Operations}

IEEE~754는 무엇보다 덧셈, 뺄셈, 곱셈, 나눗셈, 제곱근, 그리고 곱한 뒤 더하는 FMA가 \term{올바르게 반올림된다}{correctly rounded}고 약속한다. 두 수 $a,b\in\mathbb F$에 대해 컴퓨터가 내는 $a\oplus b$는 정확한 실수 $a+b$를 반올림한 $\fl(a+b)$와 같아야 하고 다른 연산도 마찬가지이다\cite{ieee754-2019}. 그래서 같은 연산은 어느 컴퓨터에서나 같은 비트를 내고 오차를 분석할 때도 연산마다 반올림 한 번만 따지면 된다. 반면 $\exp$나 $\sin$ 같은 초월 함수에 대해서는 표준이 올바른 반올림을 권할 뿐 요구하지 않는다. 올바른 반올림을 하려면 결과가 두 수의 한가운데에 아주 가까운 경우를 가려내야 하는데 그러려면 중간 계산을 얼마나 정밀하게 해야 하는지를 일반적으로는 미리 알기 어렵기 때문이다(Table Maker의 딜레마). binary64의 여러 함수에서는 최악의 경우를 이미 다 찾아 두었으므로 올바르게 반올림하는 구현도 있다. 그래서 라이브러리마다 초월 함수의 마지막 비트가 다를 수 있다(13--14장).

\begin{theorem}[반올림 오차의 기본 정리]
\label{thm:standard-model}
$x$가 실수이고 $2^{-1022}\le|x|\le\max\mathbb F$이면
\[
\fl(x)=x(1+\delta),\qquad|\delta|\le u=2^{-53}
\]
이다. 또한 $\fl(x)=x/(1+\delta')$, $|\delta'|\le u$로도 쓸 수 있다.
\end{theorem}

\begin{proof}
$2^e\le|x|<2^{e+1}$이라 하자($e\ge-1022$). 이 이진 구간의 간격은 $2^{e-52}$이므로 가장 가까운 수까지의 거리는 그 절반 이하이고 $|\fl(x)-x|\le2^{e-53}$이다. 위 끝에서 $\fl(x)=2^{e+1}$로 반올림되는 경우에도 이 부등식은 그대로이다. 따라서 $|\delta|=|\fl(x)-x|/|x|\le2^{e-53}/2^e=2^{-53}$이다. 둘째 꼴은 $\delta'=(x-\fl(x))/\fl(x)$로 두면 되고 $|\fl(x)|\ge2^e$이므로 같은 계산으로 $|\delta'|\le2^{-53}$이다.
\end{proof}

연산이 올바르게 반올림되므로 이 정리는 곧 연산의 오차가 된다. 넘침과 언더플로가 없으면 $a\oplus b=(a+b)(1+\delta)$, $|\delta|\le u$이고 다른 연산도 같은 꼴이다. 이것을 \term{표준 모형}{standard model}이라 한다. 11장의 반올림 오차 해석은 모두 이 한 줄에서 출발한다. 경계도 따져 두자. 상대오차는 $x$가 이진 구간의 아래 끝 바로 위에 있으면서 이웃한 두 수의 한가운데에 있을 때 $u$에 한없이 가까워진다(정확한 상한은 $u/(1+u)$이다). 위 끝 근처에서는 절반 정도로 줄어든다. \cref{fig:relerr}는 $[1/4,8.25)$의 무작위 실수 $3000$개를 정확한 유리수로 만들어 반올림한 상대오차이다. 한계는 이진 구간마다 아래 끝에서 $u$에 다가가고 위 끝으로 가며 $u/2$로 줄어드는 톱니 모양이다. 모든 점이 이 한계 안에 있다. 가장 큰 값은 $0.978u$였다.

\begin{figure}[htbp]
\centering
\begin{tikzpicture}
\begin{axis}[sdplot,xlabel={$x$},ylabel={$|\fl(x)-x|/(xu)$},xmin=0.25,xmax=8.25,ymin=0,ymax=1.05,width=0.7\textwidth,height=0.4\textwidth]
\addplot[black!45,only marks,mark=*,mark size=0.4pt] table[x=x,y=rel]{figures/10/relerr.dat};
\addlegendentry{무작위 실수}
\addplot[black,thick,domain=0.25:0.5,samples=40] {0.25/x};
\addlegendentry{$2^e/x$ (한계)}
\addplot[black,thick,domain=0.5:1,samples=40,forget plot] {0.5/x};
\addplot[black,thick,domain=1:2,samples=40,forget plot] {1/x};
\addplot[black,thick,domain=2:4,samples=40,forget plot] {2/x};
\addplot[black,thick,domain=4:8,samples=40,forget plot] {4/x};
\addplot[black,thick,domain=8:8.25,samples=10,forget plot] {8/x};
\end{axis}
\end{tikzpicture}
\caption{무작위 실수를 binary64로 반올림한 상대오차를 $u$로 나눈 값(정확한 유리수 계산; \code{figures/10/fp.py}). 한계 $2^e/x$는 이진 구간 $[2^e,2^{e+1})$의 아래 끝에서 $1$, 위 끝에서 $1/2$이다.}
\label{fig:relerr}
\end{figure}

서막의 예로 돌아가자. $0.1$의 근삿값의 상대오차는 정확히 $2^{-54}=u/2\approx5.55\times10^{-17}$이다. \code{0.1}과 \code{0.2}의 근삿값을 정확히 더하면 $10808639105689191/2^{55}=0.30000000000000001665\cdots$이다. $[1/4,1/2)$의 간격 $2^{-54}$로 재면 이 값의 이웃은 \code{0.3}의 근삿값 $0.29999999999999998889\cdots$와 그다음 수 $0.30000000000000004440\cdots$이다. 정확한 합은 두 수의 정확히 한가운데에 있고(정확한 유리수로 확인했다) 짝수로 가르는 규칙이 뒤의 수를 고른다. 비트 패턴으로 보면 두 수는 \code{0x3FD3333333333333}과 \code{0x3FD3333333333334}이다. 마지막 비트가 $0$인 쪽, 곧 유효숫자가 짝수인 쪽이 뒤의 수이다(\cref{fig:point-three}).

\begin{figure}[htbp]
\centering
\begin{tikzpicture}[x=0.72cm,y=1cm]
\draw[->] (-0.5,0) -- (11,0);
\draw[thick] (0,-0.25)--(0,0.25);
\draw[thick] (10,-0.25)--(10,0.25);
\node[anchor=north,font=\small,align=center] at (0,-0.3) {\code{0x3FD3333333333333}\\$0.29999999999999998889\cdots$};
\node[anchor=north,font=\small,align=center] at (10,-0.3) {\code{0x3FD3333333333334}\\$0.30000000000000004440\cdots$};
\draw (2,-0.12)--(2,0.12);
\node[anchor=south,font=\small] at (2,0.15) {$0.3$};
\fill (5,0) circle (2.2pt);
\draw[black!50] (5,-0.1)--(5,-1.35);
\node[anchor=north,font=\small,align=center] at (5,-1.4) {$\code{0.1}+\code{0.2}$의 정확한 합 $0.30000000000000001665\cdots$};
\draw[->,thick] (5.15,0.15) to[bend left=25] node[midway,above,font=\small] {짝수 쪽으로} (9.85,0.15);
\end{tikzpicture}
\caption{\code{0.1 + 0.2}의 반올림. 두 근삿값의 정확한 합은 이웃한 두 binary64 수(간격 $2^{-54}$)의 정확히 한가운데에 있고 짝수로의 반올림이 오른쪽 수를 고른다. 리터럴 \code{0.3}은 왼쪽 수를 가리킨다(정확한 유리수 계산).}
\label{fig:point-three}
\end{figure}


\subsection{Units in the Last Place}

반올림 오차를 상대오차 대신 ``마지막 자리 몇 개''로 재는 것이 편할 때가 많다.

\begin{definition}[ULP]
\label{def:ulp}
$2^e\le|x|<2^{e+1}$인 실수 $x$에 대해 $\operatorname{ulp}(x)=2^{\max(e,-1022)-52}$로 정의하고 $\operatorname{ulp}(0)=2^{-1074}$로 둔다.
\end{definition}

$\operatorname{ulp}(x)$는 $x$가 속한 이진 구간의 간격이다. 그러니 \cref{thm:standard-model}의 증명에서 보인 것은 $|\fl(x)-x|\le\operatorname{ulp}(x)/2$이다. ULP의 정의는 문헌마다 조금씩 다르다. 차이는 $x$가 $2$의 거듭제곱에 걸릴 때 어느 쪽 간격을 쓰느냐에 있다. 여기의 정의는 Muller 등이 Goldberg의 정의를 실수로 넓힌 것이라 부르는 정의이다\cite{goldberg1991,muller2018handbook}. 근삿값 $\hat y$와 참값 $y$의 오차를 $|\hat y-y|/\operatorname{ulp}(y)$로 재면 흔들림의 영향이 적고 정수 근처의 값이 나오므로 해석이 쉽다. ``$1$ ULP 이내''는 대략 상대오차 $2^{-52}$에서 $2^{-53}$ 이내라는 뜻이다.

두 binary64 수의 거리는 더 단순하게 잴 수 있다. \cref{prop:fp-bits-monotone}의 $\mu$에 부호를 붙인 순서 보존 사상 $\kappa(x)=\pm\mu(|x|)$를 쓰면 된다. 1장의 정렬 키와 달리 이 $\kappa$는 $-0$과 $+0$을 같은 $0$으로 보낸다. 그러면 $a$에서 $b$까지 이웃한 수로 건너가는 걸음의 수, 곧 $(a,b]$에 있는 binary64 수의 개수($\pm0$은 하나로 센다)가 $|\kappa(a)-\kappa(b)|$이다(\cref{alg:ulp-distance}). stardust의 정확도 계약은 거의 모두 이 단위로 적힌다(``Arb로 구한 참값을 binary64로 반올림한 값에서 $k$ ULP 이내'').

\begin{algorithm}[htbp]
\begin{algorithmic}[1]
\Require binary64 수 $a$, $b$
\Ensure $(\min,\max]$에 있는 binary64 수의 개수 (부호 없는 $64$비트 정수)
\State \textbf{if} 둘 다 NaN \textbf{then} \Return $0$; \textbf{if} 하나만 NaN \textbf{then} \Return $2^{64}-1$ \Comment{표지값}
\Function{Ordered}{$x$} \Comment{$\kappa$: 순서를 지키는 정수}
  \State $i\gets x$의 비트를 부호 있는 $64$비트로; \Return $i<0$ \textbf{?} $-(i\wedge\mathtt{0x7FF\ldots F})$ \textbf{:} $i$
\EndFunction
\State $h\gets\max(\Call{Ordered}{a},\Call{Ordered}{b})$; $\ell\gets\min(\cdots)$
\State \Return $h-\ell$를 부호 없는 $64$비트로 \Comment{$2^{64}$를 법으로 감기지만 참값이 $2^{64}$ 미만}
\end{algorithmic}
\caption{두 binary64 수 사이의 ULP 거리 (stardust \code{bench/accuracy/common.py}의 \code{ulp\_distance})}
\label{alg:ulp-distance}
\end{algorithm}

\cref{alg:ulp-distance}의 \code{Ordered}는 음수의 크기 부분에 마이너스를 붙여 $-0$과 $+0$을 같은 $0$으로 보낸다. 그래서 두 영의 거리는 $0$이고 $0$을 사이에 둔 두 수의 거리는 양쪽의 서브노멀 수들까지 모두 센다. $-\infty$에서 $+\infty$까지의 거리 $2^{64}-2^{53}$은 부호 있는 $64$비트 정수에 들어가지 않으므로 마지막 뺄셈을 부호 없는 정수로 한다. 2장에서 보았듯이 이 뺄셈은 $2^{64}$를 법으로 감기지만 참값이 $2^{64}$보다 작으므로 결과는 정확하다. 참값이 $0$에 가까우면 상대오차처럼 ULP 거리도 커진다. 부호가 바뀌면 $0$ 양쪽의 서브노멀 수까지 모두 세므로 $2^{62}$ 단위의 값이 나오기도 한다. 그래서 결과가 $0$에 가까운 함수의 계약은 절대오차나 $|\hat y-y|/\max(1,|y|)$ 같은 섞인 오차로 따로 적는다(28장). 복소수 결과는 실수부와 허수부의 오차 가운데 큰 쪽을 참값의 두 부분 가운데 큰 쪽의 간격으로 나누어 잰다(\code{bench/accuracy/common.py}의 \code{ulps}). 그러면 한 부분이 $0$에 가깝더라도 그 근처의 수많은 작은 수까지 세느라 거리가 부풀지 않는다.

ULP 거리와 상대오차의 관계는 흔들림만큼 느슨하다.

\begin{proposition}
\label{prop:ulp-relerr}
같은 부호의 정규 binary64 수 $a$, $b$의 ULP 거리를 $d\ge1$이라 하고 $m=\min(|a|,|b|)$, $M=\max(|a|,|b|)$라 하면
\[
d\cdot2^{-53}m<|a-b|\le d\cdot2^{-52}M
\]
이다. 따라서 $\hat y$가 $\fl(y)$에서 $d$ ULP 떨어져 있으면 $\hat y$의 상대오차는 대략 $(d-1)\cdot2^{-53}$보다 크고 $(d+1/2)\cdot2^{-52}$보다 작다.
\end{proposition}

\begin{proof}
$a$에서 $b$까지의 $d$걸음은 모두 $[m,M]$ 안의 이웃한 수 사이를 건너고 한 걸음의 크기는 그 자리의 이진 구간 $[2^e,2^{e+1})$의 간격 $2^{e-52}$이다. 걸음이 시작하는 수 $z$는 $2^e\le|z|<2^{e+1}$이므로 $2^{-53}|z|<2^{e-52}\le2^{-52}|z|$이고 $m\le|z|\le M$이다. $d$개를 더하면 부등식이 나온다. 둘째 주장은 $|\fl(y)-y|\le u|y|$와 $m$, $M$이 $|y|(1+O(d\,2^{-52}))$라는 데서 나온다.
\end{proof}

그래서 ``$k$ ULP 이내''라는 계약은 상대오차로 옮기면 이진 구간 안의 위치에 따라 두 배까지 차이가 난다. 거꾸로 상대오차 $c\cdot2^{-52}$ 이내는 대략 $2c+1$ ULP 이내이다.

\subsection{Exact Operations}

결과가 이미 $\mathbb F$에 있으면 반올림은 오차를 만들지 않고 연산은 정확하다. $2$의 거듭제곱을 곱하거나 나누면 넘침과 언더플로가 없는 한 지수만 바뀌므로 정확하다. \cref{prop:fp-integers}의 범위에서 정수끼리 더하고 곱하는 것도 결과가 $2^{53}$ 이하이면 정확하다. 이보다 덜 자명한 정확성의 원천이 가까운 두 수의 뺄셈이다.

\begin{lemma}[Sterbenz]
\label{lem:sterbenz}
$x,y\in\mathbb F$가 $y/2\le x\le2y$를 만족하면 $x-y\in\mathbb F$이다. 따라서 $x\ominus y=x-y$는 정확하다.
\end{lemma}

\begin{proof}
$x$와 $y$는 같은 부호이므로 둘 다 양수라 해도 되고 대칭성에서 $y\le x\le2y$라 하자. $y$의 지수를 $e$라 하면 $x$와 $y$는 모두 $2^{\max(e,-1022)-52}$의 정수배이다. $y$는 정의에서, $x$는 $x\ge y$이므로 지수가 $e$ 이상이어서 그 간격이 이 수의 배수이기 때문이다. 따라서 $x-y$도 이 수의 정수배이다. 한편 $0\le x-y\le y<2^{e+1}$이므로 $x-y$는 이 간격의 $2^{53}$배 미만이고 유효숫자 $53$비트로 적힌다. 그러므로 $x-y\in\mathbb F$이다.
\end{proof}

가까운 수의 뺄셈은 상대오차를 크게 키우는 \textbf{소거}의 원인으로 악명이 높다. 그러나 Sterbenz의 보조정리에 따르면 소거는 뺄셈 자체에서는 오차를 만들지 않는다. 소거가 위험한 것은 뺄셈이 \textbf{이미 오차를 품은} 두 값의 앞자리를 지워 그 오차를 드러내기 때문이고 11장에서 이 구별을 정식으로 다룬다. 정확한 뺄셈은 그 자체로 쓸모가 있다. 12장의 오차 없는 변환과 13장의 인자 축소에서 거듭 쓰이고 이 책에서도 이미 몇 번 나왔다. 예컨대 stardust의 NEON \code{expm1}(\code{vmath.hpp})은 $e^x=s(1+t)$, $s=2^{k+j/128}$로 나눈 뒤 $e^x-1=(s-1)+st$를 $\mathrm{fma}(s,t,s-1)$로 계산한다. $x$가 $-\log2$와 $\log2$ 사이에 있으면 $s\in[1/2,2]$이므로 \cref{lem:sterbenz}에 따라 $s-1$이 정확하고 반올림은 마지막 FMA 한 번뿐이다. 다른 범위의 처리는 13장에서 다룬다. $2$의 거듭제곱을 곱하면 정확하다는 성질도 같은 식으로 쓴다. stardust의 복소 절댓값은 큰 부분이 $2^{500}$보다 크거나 $2^{-500}$보다 작으면 두 부분에 $2^{\mp600}$을 곱해 크기를 맞춘다. 그 뒤 $\sqrt{a^2+b^2}$를 계산하고 다시 $2^{\pm600}$을 곱한다. 큰 부분을 곱한 값은 정확하고 작은 부분이 이때 언더플로해도 그 크기가 결과의 마지막 자리에 닿지 못한다(\code{scalar\_ops.hpp}의 \code{complex\_abs}).

Sterbenz의 조건은 충분조건일 뿐이다. \cref{fig:sterbenz}에서 $[1/4,4]^2$의 무작위 쌍 $4000$개 가운데 쐐기 $y/2\le x\le2y$ 안의 $1755$쌍은 모두 정확했고 쐐기 밖의 $2245$쌍 가운데서도 $1038$쌍이 정확했다. 밖의 쌍은 차의 유효숫자가 $53$비트에 들어갈 때만, 곧 두 수의 지수 차와 끝자리의 $0$에 따라 우연히 정확하다.

\begin{figure}[htbp]
\centering
\begin{tikzpicture}
\begin{loglogaxis}[sdplot,xlabel={$x$},ylabel={$y$},xmin=0.25,xmax=4,ymin=0.25,ymax=4,width=0.5\textwidth,height=0.5\textwidth,log basis x=2,log basis y=2]
\addplot[black!25,only marks,mark=*,mark size=0.5pt] table[x=x,y=y,restrict expr to domain={\thisrow{exact}}{0:0}]{figures/10/sterbenz.dat};
\addlegendentry{부정확}
\addplot[black,only marks,mark=*,mark size=0.5pt] table[x=x,y=y,restrict expr to domain={\thisrow{exact}}{1:1}]{figures/10/sterbenz.dat};
\addlegendentry{정확}
\addplot[black,thick,domain=0.25:2] {2*x};
\addlegendentry{$y=2x$, $y=x/2$}
\addplot[black,thick,domain=0.5:4,forget plot] {x/2};
\end{loglogaxis}
\end{tikzpicture}
\caption{$x\ominus y$가 정확한 쌍(검정)과 부정확한 쌍(회색). 두 직선 사이의 쐐기가 \cref{lem:sterbenz}의 영역이고 그 안에서는 예외가 없다(정확한 유리수로 판정; \code{figures/10/sterbenz.py}).}
\label{fig:sterbenz}
\end{figure}

\subsection{Double Rounding}

반올림을 두 번 하면 한 번 한 것과 다를 수 있다. 정확한 값을 먼저 더 넓은 형식으로 반올림하고 다시 binary64로 반올림한다고 하자. 첫 반올림이 binary64의 한가운데 값을 만들면 둘째 반올림이 잘못된 쪽을 고를 수 있다. 예를 들어 정확한 값이 $1+2^{-53}+2^{-80}$이면 binary64로는 위로 가야 한다($1+2^{-52}$). 그러나 유효숫자 $64$비트인 x87의 확장 형식으로 먼저 반올림하면 $2^{-80}$이 사라져 정확히 한가운데 $1+2^{-53}$이 되고 다시 짝수로 반올림하면 $1$이 된다. 다행히 넓은 형식이 충분히 넓으면 이런 일이 기본 연산에서는 일어나지 않는다.

\begin{theorem}[Figueroa]
\label{thm:double-rounding}
두 반올림이 모두 가장 가까운 수로의 반올림이고 넘침과 언더플로가 없다고 하자. 정밀도 $p$의 두 수에 대한 $+$, $-$, $\times$, $\div$, $\sqrt{\ }$의 결과를 정밀도 $q\ge2p+2$로 먼저 반올림하고 다시 정밀도 $p$로 반올림한 값은 한 번에 정밀도 $p$로 반올림한 값과 같다.
\end{theorem}

곱셈은 바로 보인다. 정밀도 $p$인 두 유효숫자의 곱은 $2p$비트 이하의 정수이므로 $q\ge2p$이면 첫 반올림이 아무것도 바꾸지 않는다. 나머지 연산의 줄거리는 이렇다. 정확한 결과는 정밀도 $p$ 격자의 한가운데에 정확히 놓이거나 그렇지 않으면 한가운데에서 정밀도 $2p$ 남짓한 격자의 한 칸 이상 떨어져 있다. 덧셈은 두 수의 지수 차가 크면 작은 수가 한가운데에서 떨어진 거리를 정하는데 그 거리가 아래쪽 자리들에 흩어져 있다. 나눗셈 $a/b$는 나머지 $a-qb$가 정수 격자 위에 있고 제곱근 $\sqrt a$는 $a-r^2$이 그렇다. 그래서 $q$가 충분히 크면 첫 반올림이 한가운데가 아닌 값을 한가운데로 옮길 수 없다(Figueroa\cite{figueroa1995}, 연산별 증명은 \cite{muller2018handbook}). 정밀도 $4$와 $5$에서 모든 쌍을 검사하면 연산별로 필요한 $q$는 곱셈과 나눗셈이 $2p$, 덧셈과 뺄셈이 $2p+1$, 제곱근이 $2p+2$이다(\cref{tab:double-rounding}). 그러니 정리의 $2p+2$는 제곱근에서 나온 값이다. 무작위 표본으로는 이 경계를 볼 수 없다. 첫 반올림이 일을 그르치는 비율은 대략 $2^{-(q-p)-1}$로 줄어들어 $q=70$쯤이면 $2\times10^5$개 표본에서 한 번도 나오지 않는다. 경계 근처의 그런 드문 경우는 정리가 잡아 준다.

\begin{table}[htbp]
\begin{sdtable}{colspec={lrrrrr}, row{1}={font=\small\bfseries}}
\toprule
binary64로 두 번 반올림 & $q=54$ & $56$ & $60$ & $64$ & $70$ \\
\midrule
무작위 곱 $2\times10^5$개 & $0.249$ & $0.062$ & $3.8\times10^{-3}$ & $2.4\times10^{-4}$ & $0$ \\
무작위 합 $2\times10^5$개 & $0.161$ & -- & $2.5\times10^{-3}$ & $1.2\times10^{-4}$ & $0$ \\
\midrule
전수 검사에서 안전한 최소 $q$ & $+$, $-$ & $\times$ & $\div$ & $\sqrt{\ }$ & \\
$p=4$ & $9$ & $8$ & $8$ & $10$ & \\
$p=5$ & $11$ & $10$ & $10$ & $12$ & \\
\bottomrule
\end{sdtable}
\caption{이중 반올림. 위는 정밀도 $q$를 거쳐 binary64로 반올림했을 때 한 번 반올림한 값과 다른 비율(합은 두 수의 지수 차를 $0$--$80$에서 뽑았다), 아래는 작은 정밀도의 모든 쌍을 검사해 얻은 안전한 최소 $q$이다(정확한 유리수 계산; \code{figures/10/double\_round\_*.py}).}
\label{tab:double-rounding}
\end{table}
 binary64에서 이 조건은 $q\ge108$이다. binary128($q=113$)은 충분하지만 x87의 $64$비트는 모자라고 위 예가 그 경우이다. binary32의 연산을 binary64로 계산하고 반올림하는 것은 $53\ge2\cdot24+2$이므로 안전하다. stardust는 binary64 하나만 쓰고 확장 형식을 거치지 않으므로(AArch64에는 binary64보다 넓은 부동소수점 산술 형식이 없고 \code{long double}도 binary64이다) 기본 연산에서 이 문제를 만나지 않는다. 다만 double-double 같은 두 단계 계산(12장)이나 십진수 변환(15장)에서는 같은 함정을 따로 피해야 한다.
\section{Subnormal Numbers and Gradual Underflow}
\label{sec:subnormal}

정규 수만 있다면 $0$과 가장 작은 정규 수 $2^{-1022}$ 사이에 큰 빈틈이 생긴다. 그 빈틈의 너비는 $2^{-1022}$이지만 바로 위의 간격은 $2^{-1074}$이므로 빈틈이 이웃한 간격보다 $2^{52}$배나 넓다. 계산 결과가 이 빈틈에 떨어지면 $0$으로 보내는 수밖에 없고 이것을 \term{갑작스런 언더플로}{abrupt underflow} 또는 영으로의 플러시라 한다. 서브노멀 수는 이 빈틈을 정규 수의 마지막 간격 $2^{-1074}$로 고르게 채운다. 그러면 값이 작아질 때 정밀도가 한꺼번에 사라지지 않고 비트 하나씩 줄어들게 한다(\cref{fig:subnormal}). 이것이 \term{점진적 언더플로}{gradual underflow}이다.

\begin{figure}[htbp]
\centering
\begin{tikzpicture}[x=0.48cm,y=1cm]
\draw[->] (0,1.6) -- (21,1.6);
\foreach \i in {0,...,8} \draw ({4+\i*0.5},1.48)--({4+\i*0.5},1.72);
\foreach \i in {0,...,8} \draw ({8+\i},1.48)--({8+\i},1.72);
\draw[thick] (0,1.4)--(0,1.8) node[above] {$0$};
\draw[thick] (4,1.4)--(4,1.8) node[above] {$2^{e_{\min}}$};
\draw[thick] (8,1.4)--(8,1.8) node[above] {$2^{e_{\min}+1}$};
\draw[thick] (16,1.4)--(16,1.8) node[above] {$2^{e_{\min}+2}$};
\node[anchor=west,font=\small] at (21.2,1.6) {플러시};
\draw[->] (0,0) -- (21,0);
\foreach \i in {0,...,7} \draw ({\i*0.5},-0.12)--({\i*0.5},0.12);
\foreach \i in {0,...,8} \draw ({4+\i*0.5},-0.12)--({4+\i*0.5},0.12);
\foreach \i in {0,...,8} \draw ({8+\i},-0.12)--({8+\i},0.12);
\draw[thick] (0,-0.2)--(0,0.2);
\draw[thick] (4,-0.2)--(4,0.2);
\draw[thick] (8,-0.2)--(8,0.2);
\draw[thick] (16,-0.2)--(16,0.2);
\draw[decorate,decoration={brace,mirror,amplitude=4pt}] (0,-0.3)--(4,-0.3) node[midway,below=4pt,font=\small] {서브노멀};
\node[anchor=west,font=\small] at (21.2,0) {점진적};
\draw[decorate,decoration={brace,mirror,amplitude=4pt}] (0.05,1.3)--(3.95,1.3) node[midway,below=4pt,font=\small] {빈틈};
\end{tikzpicture}
\caption{정밀도 $p=4$인 장난감 형식의 양수 부분. 위는 서브노멀 수가 없을 때로 $0$과 $2^{e_{\min}}$ 사이가 비어 있고 아래는 서브노멀 수가 그 구간을 $2^{e_{\min}}$ 바로 위의 간격으로 채운 모습이다.}
\label{fig:subnormal}
\end{figure}


\begin{proposition}
\label{prop:gradual-underflow}
서브노멀 수가 있는 binary64에서 $x,y\in\mathbb F$가 유한하고 $x\ne y$이면 $x\ominus y\ne0$이다. 서브노멀 수를 $0$으로 플러시하면 이 성질이 깨진다.
\end{proposition}

\begin{proof}
모든 유한한 binary64 수는 $2^{-1074}$의 정수배이므로 $x-y$도 $0$이 아닌 $2^{-1074}$의 정수배이다. $|x-y|<2^{-1022}$이면 $x-y=k\cdot2^{-1074}$, $0<|k|<2^{52}$이므로 $x-y$는 서브노멀 수로 정확히 나타나고 $x\ominus y=x-y\ne0$이다. $|x-y|\ge2^{-1022}$이면 반올림 결과의 크기도 $2^{-1022}$ 이상이다. 플러시하는 경우의 반례로 $x=1.5\cdot2^{-1022}$, $y=2^{-1022}$를 들 수 있다. 두 수는 모두 정규 수이지만 차 $2^{-1023}$은 서브노멀 수이므로 $0$이 된다.
\end{proof}

이 성질이 없으면 \code{if (x != y) z = 1 / (x - y);} 같은 코드가 $0$으로 나눌 수 있다. 점진적 언더플로가 있으면 표준 모형도 언더플로 영역까지 넓어진다. 반올림 결과가 서브노멀 영역에 떨어지면 상대오차는 $u$로 묶이지 않는다. 그래도 절대오차는 간격의 절반 $2^{-1075}$로 묶이므로 모든 유한한 결과에 대해
\begin{equation}
\fl(x)=x(1+\delta)+\eta,\qquad|\delta|\le u,\quad|\eta|\le2^{-1075},\quad\delta\eta=0
\label{eq:standard-model-underflow}
\end{equation}
이 성립한다. 곱셈과 나눗셈에는 이 꼴이 필요하다. 덧셈과 뺄셈은 \cref{prop:gradual-underflow}의 증명에서 보았듯이 결과가 서브노멀 영역에 떨어지면 언제나 정확하므로 $\eta=0$으로 둘 수 있다. 서브노멀 영역에서 정밀도가 줄어드는 모습은 측정으로도 보인다. stardust에서 \code{3e-310 * 0.5}는 $1.49999999999997\times10^{-310}$으로 유효숫자가 열넷 남짓만 남고 가장 작은 서브노멀 수의 절반 \code{5e-324 * 0.5}는 정확히 한가운데여서 짝수 쪽인 $0$으로 간다.

1970년대의 컴퓨터는 회사마다 부동소수점 형식과 반올림 규칙이 달라서 같은 프로그램이 기계를 옮기면 다른 답을 내는 일이 흔했다. 1977년에 시작된 IEEE의 표준화 위원회는 점진적 언더플로를 두고 가장 오래 논쟁했다. 다른 회사들은 서브노멀 수가 하드웨어를 복잡하고 느리게 만든다며 반대했고 Kahan 쪽은 \cref{prop:gradual-underflow} 같은 성질이 프로그램을 증명할 수 있게 해 준다고 맞섰다\cite{coonen1981underflow}. 1985년의 표준은 Kahan의 안을 거의 그대로 받아들였고 Kahan은 이 공로로 1989년 Turing상을 받았다\cite{severance1998}.

느리다는 반대는 오늘날에도 일부 맞는다. 많은 x86 프로세서는 서브노멀 입력이나 결과가 나오면 마이크로코드의 도움을 받아 처리하므로 그 연산이 정규 수의 연산보다 수십 배 느려질 수 있고 그래서 성능을 위해 플러시 모드(x86의 FTZ와 DAZ, ARM의 FPCR.FZ)를 켜는 프로그램도 있다. 이 모드는 \cref{prop:gradual-underflow}를 깨뜨리고 같은 식의 비트를 바꾼다. stardust가 도는 Apple M4는 서브노멀을 하드웨어에서 그대로 처리한다. L1 캐시에 들어가는 $16384$개 원소를 한 스레드로 곱하면 정규 수가 $3.08\,\mu$s, 서브노멀 수가 $3.04\,\mu$s로 차이가 없었다(복제본, \code{figures/10/review\_sd.py}). 큰 배열에서는 메모리 대역폭이 시간을 정하므로 몇 사이클의 벌칙이 가려질 수 있다. 근거는 캐시 안에서 잰 값으로 삼았다($10^7$개에서도 $0.68$ms 대 $0.66$ms였다). 그래서 stardust는 플러시 모드를 켜지 않고 표준의 점진적 언더플로를 그대로 쓰며 \cref{eq:standard-model-underflow}가 모든 기본 연산에서 성립한다.

\section{Special Values}
\label{sec:special-values}

\subsection{Signed Zero}

binary64에는 $+0$과 $-0$이 따로 있다. 두 영은 \code{==}로 견주면 같지만 몇몇 연산에서는 다르게 행동한다. $1/(+0)=+\infty$이고 $1/(-0)=-\infty$이며 $\sqrt{-0}=-0$이다. 덧셈에서는 $(-0)+(-0)=-0$이다. 그러나 짝수로의 반올림에서 크기가 같고 부호가 다른 두 수의 합은 $+0$이므로 $x-x=+0$이고 $(-0)+(+0)=+0$이다. 따라서 $-0$을 더해도 모든 $x$에서 $x$가 그대로 나오지만 $+0$을 더하면 $-0$이 $+0$으로 바뀐다. 덧셈의 항등원은 $-0$이다.

이 사실은 합을 계산하는 코드에 바로 드러난다. 누적 변수를 $+0$에서 시작하면 $-0$ 하나짜리 배열의 합이 $+0$이 되고 $-0$에서 시작해야 $-0$이 나온다. stardust의 \code{sum([-0.0])}은 $+0$을 낸다. numpy도 같은 값을 낸다(복제본의 numpy로 확인). 리덕션은 누적을 $+0$에서 시작한다고 보면 된다. $+0$은 덧셈의 항등원이 아니지만 빈 합의 값이다. 그래서 \code{sum([-0.0, -0.0])}도 $+0$이다. 반면 누적합 \code{cumulative\_sum}은 첫 원소에서 시작하므로 $[-0,-0]$을 그대로 낸다. 곱 \code{prod([-0.0])}은 $1\cdot(-0)=-0$이다. \code{abs(-0.0)}은 부호 비트를 지워 $+0$이고 \code{sign(-0.0)}은 상수 $0$을 돌려주므로 $+0$이다. 결과의 부호까지 따지는 코드는 합의 부호에 기대지 말고 따로 다루어야 한다.

부호 있는 영은 $0$으로 다가간 방향의 정보를 남긴다. 아래로 넘쳐 $0$이 된 음수 $-10^{-400}$은 $-0$이 되므로 그 역수가 $-\infty$로 올바른 부호를 갖는다. 5장에서 본 복소함수의 가지 자름 위의 값도 부호 있는 영으로 구별되는데 예를 들어 $\sqrt{-1+0i}=i$이고 $\sqrt{-1-0i}=-i$로 두면 자름의 위아래에서 각각의 연속성이 지켜진다\cite{kahan1987branch}. 그래서 부호 있는 영을 다룰 때는 조심해야 한다. \code{x == 0}으로는 부호를 알 수 없으므로 부호는 \code{signbit}이나 \code{copysign}으로 읽는다.

\subsection{Infinities}

$\pm\infty$는 넘침의 결과이자 $1/0$처럼 극한이 정해진 연산의 결과이다. 무한대가 낀 산술은 극한의 규칙을 따라 $\infty+1=\infty$, $2\cdot\infty=\infty$, $1/\infty=0$이다. 비교에서는 모든 유한한 수보다 크거나 작다. 무한대가 있으니 넘침이 일어나도 계산은 멈추지 않고 이어진다. 넘친 중간값이 나중에 $1/\infty=0$처럼 의미 있는 값으로 돌아오기도 한다. 다만 그렇게 돌아온 값이 참값과 맞는다는 보장은 없다. $\sqrt{x^2+y^2}$를 그대로 계산하면 $x=10^{200}$에서 $x^2$이 넘쳐 $\infty$가 나오는데 참값은 $10^{200}$ 근처로 아무 문제가 없다. 이런 식은 크기를 먼저 맞추어 넘침을 피해야 한다(\code{hypot}, 14장). 결과에 $\infty$가 나왔다면 계산이 망가졌다는 신호일 수 있다.

\subsection{NaN}

$0/0$, $\infty-\infty$, $0\cdot\infty$, $\sqrt{-1}$처럼 극한조차 정해지지 않는 연산의 결과는 NaN(Not a Number)이다. 비트로는 $E=2047$이고 $f\ne0$인 패턴이 모두 NaN이므로 NaN은 $2^{53}-2$개나 된다. 가수 필드의 맨 위 비트가 $1$이면 \term{조용한 NaN}{quiet NaN}, $0$이면 \term{신호 NaN}{signaling NaN}이다(2008년 판부터 표준이 권하는 부호화이다). 조용한 NaN은 연산에 들어가면 예외 없이 NaN을 내고 신호 NaN은 연산에 들어가는 순간 무효 연산 예외를 일으킨다. 나머지 비트는 \term{페이로드}{payload}라 하여 NaN이 생긴 원인 같은 정보를 담는 데 쓸 수 있다. 신호 NaN도 복사, 부호 바꾸기, 절댓값, \code{copysign}에서는 예외를 일으키지 않는다(\cref{sec:bit-ops}).

표준은 NaN 둘이 들어간 연산의 결과가 어느 쪽 페이로드를 가질지 정하지 않는다\cite{ieee754-2019}. 하드웨어는 정한다. AArch64는 기본 NaN 모드(FPCR.DN)가 꺼져 있으면 첫 NaN 피연산자의 페이로드를 그대로 두고 신호 NaN이면 조용한 비트만 켠다. NaN을 새로 만들 때는 기본 NaN \code{0x7FF8000000000000}을 낸다. macOS에서 stardust로 잰 \cref{tab:nan-bits}를 보면 이 모드가 꺼져 있다. 라이브러리 함수는 사정이 또 다르다. \code{sd.log(-1)}은 Apple vForce가 정한 페이로드 \code{0x33}을 담은 NaN을 내고 numpy(libm)는 기본 NaN을 낸다. 같은 입력인데 두 라이브러리의 NaN 비트가 다르다. NaN에도 부호 비트가 있어서 \code{negative(nan)}의 \code{signbit}는 참이다. 그러니 페이로드나 NaN의 부호에 뜻을 싣는 코드는 이식할 수 없다고 보아야 한다.

\begin{table}[htbp]
\begin{sdtable}{colspec={ll}, row{1}={font=\small\bfseries}}
\toprule
식 & 결과의 비트 \\
\midrule
\code{add}($\mathrm{NaN}_a$, $\mathrm{NaN}_b$), 페이로드 \code{aaa}, \code{bbb} & \code{0x7FF8000000000aaa} \\
\code{add}($\mathrm{NaN}_b$, $\mathrm{NaN}_a$) & \code{0x7FF8000000000bbb} \\
\code{add}(신호 NaN \code{0x7FF0000000000001}, $0$) & \code{0x7FF8000000000001} \\
\code{divide(0, 0)}, \code{subtract(inf, inf)}, \code{sqrt(-1)} & \code{0x7FF8000000000000} \\
\code{sd.log(-1)} / numpy \code{log(-1)} & \code{0x7FF8000000000033} / \code{0x7FF8000000000000} \\
\code{negative(nan)} & \code{0xFFF8000000000000} \\
\bottomrule
\end{sdtable}
\caption{NaN의 비트가 전해지는 모습(복제본, \code{figures/10/review\_sd.py}). 첫 NaN 피연산자의 페이로드가 남고 신호 NaN은 조용한 비트가 켜진다.}
\label{tab:nan-bits}
\end{table}

1장에서 보았듯이 NaN은 무엇보다 어떤 값과도 순서가 없다. \code{x < y}, \code{x == y}, \code{x > y}가 모두 거짓이므로 \code{x != x}는 $x$가 NaN일 때만 참이다. 비교에 기대는 알고리즘에는 모두 이 점이 함정이 된다. 정렬은 NaN을 어디에 둘지 따로 정해야 하고(1장, 18장, 31장) 최댓값과 최솟값의 의미도 다시 정해야 한다.

\subsection{Maximum and Minimum}

C의 \code{a > b ? a : b}는 NaN이 섞이면 인자의 순서에 따라 결과가 달라진다. $a$가 NaN이면 비교가 거짓이라 $b$가 나오고 $b$가 NaN이면 역시 $b$, 곧 NaN이 나온다. 두 영에 대해서도 \code{+0 > -0}이 거짓이므로 결과가 인자의 순서를 따른다. 이런 함수로 배열의 최댓값을 구하면 NaN이 첫 원소인지 가운데 원소인지에 따라 답이 다르다. IEEE~754-2019는 순서에 기대지 않는 두 쌍을 정한다\cite{ieee754-2019}.
\begin{enumerate}
\item \code{maximum}과 \code{minimum}은 인자 하나라도 NaN이면 NaN을 내고 $-0<+0$으로 본다.
\item \code{maximumNumber}와 \code{minimumNumber}는 NaN을 빠진 값처럼 무시하고 다른 인자를 낸다. 둘 다 NaN이면 NaN이다.
\end{enumerate}
NaN을 하나의 값으로 보면 둘 다 교환법칙과 결합법칙을 만족하므로 배열의 최댓값을 어떤 순서와 어떤 묶음으로 구해도 같은 값이 나온다. 다만 서로 다른 NaN이 섞이면 결과로 나오는 NaN의 비트는 순서를 따른다. 그 덕분에 벡터 레지스터의 여러 갈래로 나누어 구하고 마지막에 합쳐도 된다. 거꾸로 C의 삼항 연산자로 만든 최댓값은 결합법칙이 깨지므로 벡터화하면 결과가 바뀔 수 있다.

\begin{algorithm}[htbp]
\begin{algorithmic}[1]
\Function{Maximum}{$a$, $b$} \Comment{IEEE \code{maximum}, NEON \code{FMAX}}
  \State \textbf{if} $a\ne a$ \textbf{then} \Return $a$; \textbf{if} $b\ne b$ \textbf{then} \Return $b$
  \State \textbf{if} $a=b$ \textbf{then} \Return (\Call{SignBit}{$a$} \textbf{?} $b$ \textbf{:} $a$) \Comment{$\pm0$이면 $+0$}
  \State \Return ($a>b$ \textbf{?} $a$ \textbf{:} $b$)
\EndFunction
\Function{MaximumNumber}{$a$, $b$} \Comment{\code{nanmax}, NEON \code{FMAXNM}(조용한 NaN)}
  \State \textbf{if} $a\ne a$ \textbf{then} \Return $b$; \textbf{if} $b\ne b$ \textbf{then} \Return $a$
  \State \Return \Call{Maximum}{$a$, $b$}
\EndFunction
\end{algorithmic}
\caption{NaN과 부호 있는 영을 정한 최댓값 (stardust \code{scalar\_ops.hpp}의 \code{fmax\_nan}, \code{NanMaximum})}
\label{alg:fmax-nan}
\end{algorithm}

stardust의 \code{sd.maximum}과 \code{sd.max}는 \cref{alg:fmax-nan}의 \code{Maximum}이고 \code{sd.nanmax}는 \code{MaximumNumber}이다. 조용한 NaN과 수에서는 이 정의가 ARM NEON의 \code{FMAX}와 \code{FMAXNM} 명령과 비트까지 같은 결과를 낸다. 그래서 벡터 커널은 명령 하나로 원소 두 개를 처리하고 남은 원소를 맡는 스칼라 코드도 같은 값을 낸다. 최솟값은 부호 비트로 고르는 쪽만 바뀐다($\pm0$이면 $-0$). 신호 NaN에서는 둘이 다르다. \code{FMAX}는 두 피연산자 가운데 신호 NaN을 먼저 골라 조용하게 바꾸어 낸다(아래의 \textproc{ProcessNaNs}). 반면 \cref{alg:fmax-nan}은 받은 NaN을 그대로 돌려준다. 그래서 신호 NaN이 든 배열에서는 \code{sd.maximum}의 결과가 원소의 자리에 따라 달라진다. 원소 다섯 개짜리 배열에서 벡터 경로를 탄 첫 원소는 \code{0x7FF8000000000001}, 스칼라로 처리한 마지막 원소는 \code{0x7FF0000000000001}이었다(\code{figures/10/review\_sd.py}). numpy에서 가져온 배열은 신호 NaN의 비트를 그대로 지니므로 실제로 일어날 수 있다. 두 경로를 맞추려면 스칼라 쪽의 NaN 처리를
\[
\textproc{ProcessNaNs}(a,b)=\begin{cases}\textproc{Quiet}(a)&a\text{가 신호 NaN}\\\textproc{Quiet}(b)&b\text{가 신호 NaN}\\a&a\ne a\\b&\text{그 밖}\end{cases}
\]
로 바꾸면 된다. 여기서 $\textproc{Quiet}(x)$는 비트에 $2^{51}$을 켜는 것이다. numpy의 \code{fmax}에 해당하는 공개 함수는 없고 NaN을 건너뛰는 규칙은 리덕션인 \code{nanmax}와 \code{nanmin}에서만 쓴다.

\subsection{Exceptions and Default Results}

표준은 다섯 가지 \term{예외}{exception}를 정하고 예외마다 \textbf{기본 결과}를 정해 계산을 멈추지 않고 내놓게 한다. 무효 연산은 NaN을, 영으로 나누기는 올바른 부호의 무한대를, 넘침은 무한대를, 언더플로는 서브노멀 수나 $0$을, 부정확은 반올림한 값을 낸다. 각 예외가 일어나면 상태 플래그가 세워지고 프로그램은 나중에 이 플래그를 확인할 수 있다. 하지만 고수준 언어에서 플래그를 들여다보는 일은 드물다. 실제로는 기본 결과 덕분에 계산이 끝까지 이어지고 결과에 남은 NaN이나 무한대를 보고 문제가 있었음을 알게 된다. numpy는 ufunc를 부를 때마다 플래그를 한 번 읽어 경고를 띄우고 \code{errstate}로 그 동작을 바꾸게 한다. 원소마다 읽는 것이 아니어서 비용은 작다. $10^6$개의 \code{divide(0, 0)}은 경고를 켜든 끄든 $120\,\mu$s였다(복제본). stardust는 경고를 일절 내지 않고 기본 결과만 돌려준다. 라이브러리 전체의 규칙에 따라 잘못된 인자는 예외(\code{ValueError} 등)로, 계산 중의 실패는 NaN으로 알린다. 플래그에 기대지 않는 이유는 또 있다. 스레드 풀의 일꾼 스레드에서 세워진 플래그는 부른 스레드에서 보이지 않는다. 대신 stardust는 부동소수점 환경이 기본값 그대로라고 가정한다. 반올림 모드는 짝수로의 반올림이고 플러시와 기본 NaN 모드는 꺼져 있으며 예외 트랩은 없다. 코드 어디에도 \code{fesetround}나 \code{fetestexcept}가 없다. \code{sd.round}가 \code{std::rint}로 짝수로의 반올림을 얻는 것도 이 가정에 기댄다(\cref{tab:exceptions}).

\begin{table}[htbp]
\begin{sdtable}{colspec={llll}, row{1}={font=\small\bfseries}}
\toprule
예외 & 예 & 기본 결과 & stardust의 결과 \\
\midrule
무효 연산 & $0/0$, $\infty-\infty$, $\sqrt{-1}$ & 기본 NaN & \code{0x7FF8000000000000} \\
영으로 나누기 & $1/(-0)$ & $-\infty$ & $-\infty$ \\
넘침 & $10^{308}\cdot10$ & $+\infty$ & $+\infty$ \\
언더플로 & $2^{-1074}\cdot0.5$ & $0$ 또는 서브노멀 수 & $0$ \\
부정확 & $1/3$ & 반올림한 값 & $0.3333333333333333$ \\
\bottomrule
\end{sdtable}
\caption{다섯 가지 예외와 기본 결과(복제본에서 확인). stardust는 플래그를 읽지 않고 경고도 내지 않는다.}
\label{tab:exceptions}
\end{table}

\subsection{Complex Special Values}

\code{complex128}은 binary64 두 개이므로 실수부와 허수부가 따로 이 절의 규칙을 따른다. 그러나 복소 연산의 결과는 두 부분이 섞여 나오므로 한 부분의 무한대가 다른 부분의 NaN을 낳기 쉽다. $(\infty+0i)(\infty+0i)$의 허수부는 $\infty\cdot0+0\cdot\infty$로 NaN이 되고 $(\infty+\mathrm{NaN}\,i)\cdot2$는 실수부도 $\infty\cdot2-\mathrm{NaN}\cdot0$으로 NaN이 된다. C99의 부록 G는 이런 결과에서 ``무한대인 부분이 있으면 무한대''라는 뜻을 되살리는 복구 단계를 정한다. Python의 복소수도 비슷하게 동작한다. numpy는 이 단계를 두지 않는다. stardust도 numpy에 맞추어 복구 없는 공식을 직접 적는다(\code{scalar\_ops.hpp}; \code{std::complex}의 연산자는 복구를 넣으므로 쓰지 않는다). \cref{tab:complex-special}은 stardust에서 잰 값이고 모두 numpy와 같다. 중간 넘침도 눈여겨볼 만하다. $(10^{300}+10^{300}i)^2$의 참값은 $2\cdot10^{600}i$로 실수부가 $0$이지만 $10^{600}$이 넘치므로 $-\infty+\infty i$가 나온다. \code{complex\_abs}의 $2^{\pm600}$ 곱하기가 바로 크기를 맞추어 이 넘침을 피하는 장치이다(\cref{sec:rounding}).

\begin{table}[htbp]
\begin{sdtable}{colspec={ll}, row{1}={font=\small\bfseries}}
\toprule
식 & stardust와 numpy의 결과 \\
\midrule
$(\infty+0i)\cdot(\infty+0i)$ & $\infty+\mathrm{NaN}\,i$ \\
$(\infty+\mathrm{NaN}\,i)\cdot2$ & $\mathrm{NaN}+\mathrm{NaN}\,i$ \\
$(10^{300}+10^{300}i)^2$ & $-\infty+\infty i$ \\
$(1+i)/0$ & $\infty+\infty i$ \\
$0/0$ (복소) & $\mathrm{NaN}+\mathrm{NaN}\,i$ \\
$|\infty+\mathrm{NaN}\,i|$ & $\infty$ \\
$|3\cdot10^{-310}+4\cdot10^{-310}i|$ & $5\cdot10^{-310}$ \\
\bottomrule
\end{sdtable}
\caption{복소수의 특수값(복제본, \code{figures/10/review\_sd.py}). 부록 G의 복구가 없으므로 무한대가 NaN으로 바뀌는 경우가 있다.}
\label{tab:complex-special}
\end{table}

\cref{tab:special-sd}는 이 절의 규칙을 stardust에서 직접 확인한 것이다. 모든 항목이 표준의 기본 결과와 맞는다.

\begin{table}[htbp]
\begin{sdtable}{colspec={ll|ll}, row{1}={font=\small\bfseries}}
\toprule
식 & 결과 & 식 & 결과 \\
\midrule
\code{add(-0.0, +0.0)} & $+0$ & \code{maximum(-0.0, +0.0)} & $+0$ \\
\code{add(-0.0, -0.0)} & $-0$ & \code{minimum(+0.0, -0.0)} & $-0$ \\
\code{subtract(1.0, 1.0)} & $+0$ & \code{maximum(nan, 1.0)} & NaN \\
\code{sqrt(-0.0)} & $-0$ & \code{max([1, nan, 3])} & NaN \\
\code{divide(1.0, -0.0)} & $-\infty$ & \code{nanmax([1, nan, 3])} & $3$ \\
\code{sum([-0.0])} & $+0$ & \code{nextafter(0.0, -1.0)} & $-2^{-1074}$ \\
\code{multiply(3e-310, 0.5)} & $1.49999999999997\times10^{-310}$ & \code{nextafter(max, inf)} & $+\infty$ \\
\code{cumulative\_sum([-0.0, -0.0])} & $[-0,-0]$ & \code{nextafter(inf, 0.0)} & \code{max} \\
\code{prod([-0.0])} & $-0$ & \code{sign(-0.0)}, \code{abs(-0.0)} & $+0$, $+0$ \\
\bottomrule
\end{sdtable}
\caption{stardust에서 잰 특수값의 결과(복제본, \code{figures/10/special\_sd.py}). 스칼라는 원소 하나짜리 배열로 넣었다. \code{max}는 binary64의 가장 큰 유한 수이다.}
\label{tab:special-sd}
\end{table}

\section{Conversion to and from Integers}
\label{sec:conversion}

\subsection{From Integers}

정수 $n$을 binary64로 바꾸는 것은 반올림 $\fl(n)$이다. $|n|\le2^{53}$이면 정확하고 그 위에서는 짝수로의 반올림으로 끝자리가 사라진다. stardust에서 \code{int64} 배열 $[2^{53}+1]$을 \code{float64}로 바꾸면 $9007199254740992=2^{53}$이 나온다. 이 변환은 C++의 \code{static\_cast}이고 AArch64의 \code{SCVTF} 명령 하나가 현재 반올림 모드로 올바르게 반올림한다. 명령이 없으면 \cref{alg:int-to-float}를 따르면 된다. 버리는 꼬리 $r$을 절반 $h$와 견주기만 하면 된다. 하드웨어는 이 비교를 꼬리의 맨 위 비트(가드 비트)와 그 아래 비트들의 논리합(스티키 비트)으로 한다.

\begin{algorithm}[htbp]
\begin{algorithmic}[1]
\Require 정수 $N$ (임의의 길이)
\Ensure $\fl(N)$ (짝수로의 반올림)
\State \textbf{if} $N=0$ \textbf{then} \Return $+0$
\State $\sigma\gets\operatorname{sign}N$; $M\gets|N|$; $k\gets M$의 비트 길이
\State \textbf{if} $k\le53$ \textbf{then} \Return $\sigma M$ \Comment{정확}
\State $d\gets k-53$; $q\gets M\gg d$; $r\gets M\bmod2^d$; $h\gets2^{d-1}$
\State \textbf{if} $r>h$ \textbf{or} ($r=h$ \textbf{and} $q$가 홀수) \textbf{then} $q\gets q+1$ \Comment{가드 $[r\ge h]$, 스티키 $[r\ne h]$}
\State \textbf{if} $q=2^{53}$ \textbf{then} $q\gets2^{52}$; $d\gets d+1$ \Comment{올림이 자리를 넘기면 지수를 올린다}
\State \Return $\sigma q2^d$ \Comment{비트로는 $(\sigma<0)\ll63\mid(d+1075)\ll52\mid(q-2^{52})$}
\end{algorithmic}
\caption{정수를 binary64로 올바르게 반올림하기 (넘침은 따로 검사한다)}
\label{alg:int-to-float}
\end{algorithm}

\cref{alg:int-to-float}가 맞는 까닭은 $M$이 이진 구간 $[2^{k-1},2^k)$에 있고 그 간격이 $2^d$이기 때문이다. $q2^d\le M<(q+1)2^d$이므로 두 후보가 $M$의 이웃이고 $r$과 $h$의 비교가 어느 쪽이 가까운지를 정한다. 유리수 $a/b$도 같은 방법으로 반올림할 수 있다. 나눗셈의 몫과 나머지가 $q$와 $r$ 노릇을 한다. 이 장의 그림에 쓴 정확한 반올림은 모두 이렇게 계산했다. stardust에서 $2^{62}+2^9$는 한가운데이므로 짝수 쪽 $2^{62}$로, $2^{62}+2^9+1$은 $2^{62}+2^{10}$으로 간다(복제본). 시험할 때는 Python의 \code{float(int)}가 올바르게 반올림하므로 그것과 견주면 된다. $64$비트 정수 배열을 \code{float64}로 바꾸어 계산하는 라이브러리가 흔하므로 정수의 합이나 평균이 $2^{53}$을 넘는 큰 값을 다룰 때는 이 경계를 의식해야 한다.

순서를 바꾸면 이 손실을 대부분 피할 수 있다. 원소를 하나씩 \code{float64}로 바꾸어 더하지 않고 정확한 정수합을 먼저 구한 뒤 마지막에 한 번만 반올림하면 합은 참값을 올바르게 반올림한 값이 된다. stardust의 정수 배열 \code{mean}이 이렇게 한다(\code{csrc/base/int\_sum.hpp}의 \code{IntSum}). 값마다 위 $32$비트(부호 있음)는 \code{int64} 누적기에, 아래 $32$비트는 \code{uint64} 누적기에 따로 더하고 $2^{31}$개마다 \code{\_\_int128}로 옮기므로 넘치지 않는다. 정수 덧셈은 순서와 무관하므로 어떻게 나누어 더해도 같은 비트가 나온다. 그다음 합을 한 번 반올림하고 개수로 나누니 반올림은 모두 두 번이다. $[2^{53}+1,2^{53}+1,2^{53}+1,0]$의 평균의 참값은 $6755399441055744.75$이다. 정수합 $3(2^{53}+1)$을 반올림하면 $3\cdot2^{53}+4$가 되고 $4$로 나누면 $6755399441055745$로 참값을 올바르게 반올림한 값과 같다. 원소를 먼저 바꾸면 세 원소가 모두 $2^{53}$이 되어 $6755399441055744$가 나오고 numpy의 \code{mean}이 이 값을 낸다(둘 다 복제본에서 확인). $2^{31}$개마다 옮기면 위쪽 합의 크기가 $2^{31}\cdot2^{31}=2^{62}$, 아래쪽 합이 $2^{31}(2^{32}-1)<2^{63}$을 넘지 않는다. 반올림이 두 번이므로 평균이 늘 올바르게 반올림되지는 않는다. $[2^{53}+1,0,0]$에서 합 $2^{53}+1$은 $2^{53}$으로 반올림되고 $3$으로 나누면 $3002399751580330.5$가 된다. 그런데 참값 $3002399751580331$은 binary64로 정확히 나타난다(stardust와 numpy가 같은 값을 낸다). 표준 모형으로 보면 오차는 $|\hat m-m|\le(2u+u^2)|m|$이다. 정수합의 자세한 구현은 29장에서 다룬다.

\subsection{To Integers}

반대 방향으로 binary64를 정수로 바꿀 때는 세 가지를 정해야 한다. 소수 부분을 어떻게 버릴지, 정수형의 범위를 넘는 값을 어떻게 할지, NaN을 무엇으로 할지이다. C와 C++는 소수 부분을 $0$ 쪽으로 자르도록 정하지만 결과가 정수형의 범위를 벗어나거나 값이 NaN이면 \textbf{정의되지 않은 동작}으로 둔다\cite{iso14882-2020}. 2장에서 부호 있는 정수가 넘칠 때 본 것과 같은 함정이다. 컴파일러는 그런 입력이 없다고 가정하고 코드를 바꿀 수 있다. 하드웨어도 제각각이다. x86의 \code{CVTTSD2SI}는 범위 밖과 NaN에서 모두 $-2^{63}$이라는 하나의 ``정수 무한''을 내고 ARM의 \code{FCVTZS}는 범위 밖을 가까운 끝으로 포화시키고 NaN을 $0$으로 보낸다.

\begin{algorithm}[htbp]
\begin{algorithmic}[1]
\Require binary64 수 $x$
\Ensure \code{int64} 값 (ARM \code{FCVTZS}와 같다)
\State \textbf{if} $x\ne x$ \textbf{then} \Return $0$
\State \textbf{if} $x\ge2^{63}$ \textbf{then} \Return $2^{63}-1$ \Comment{$+\infty$ 포함}
\State \textbf{if} $x<-2^{63}$ \textbf{then} \Return $-2^{63}$ \Comment{$-2^{63}$ 자신은 정확히 들어간다}
\State \Return $\operatorname{trunc}(x)$ \Comment{$|x|<2^{63}$의 정수 부분, 정확}
\end{algorithmic}
\caption{포화하는 \code{float64} $\to$ \code{int64} 변환 (stardust \code{cast.hpp}, \code{vcvtq\_s64\_f64})}
\label{alg:float-to-int}
\end{algorithm}

stardust의 \code{astype}은 \cref{alg:float-to-int}를 따른다. 경계는 $2^{63}$과 비교해도 된다. $2^{63}$은 binary64로 정확히 나타나고 그보다 작은 가장 큰 binary64 수가 $2^{63}-2^{10}$이므로 $x<2^{63}$이면 $\operatorname{trunc}(x)$가 \code{int64}에 들어가기 때문이다. 4행의 $\operatorname{trunc}$도 정확하다. $|x|\ge2^{52}$인 수는 이미 정수이고 그보다 작은 수는 정수 부분이 $53$비트 안에 들어간다. 이 네 줄을 C++로 분기 세 개를 두어 적으면 정의된 동작은 얻지만 루프가 벡터화되지 않는다. ARM의 \code{FCVTZS} 명령 하나가 정확히 이 일을 하므로 stardust는 연속 배열에서 NEON의 \code{vcvtq\_s64\_f64}를 두 번 불러 원소 네 개씩 바꾼다. 남은 원소와 간격이 있는 배열은 스칼라 \code{vcvtd\_s64\_f64}(역시 \code{FCVTZS})로 처리한다. $10^6$개 원소를 한 스레드로 바꾸면 $118\,\mu$s로 한 스레드인 numpy의 $247\,\mu$s의 절반이고 기본 스레드 설정($10$개)에서는 $28\,\mu$s이다(복제본, \code{figures/10/review\_sd.py}). 앞의 두 배가 명령의 몫이고 나머지는 스레드 풀의 몫이다. \cref{tab:astype}이 그 결과이다.

\begin{table}[htbp]
\begin{sdtable}{colspec={ll|ll}, row{1}={font=\small\bfseries}}
\toprule
입력 & \code{int64} 결과 & 입력 & \code{int64} 결과 \\
\midrule
$2.9$ & $2$ & $10^{19}$ & $2^{63}-1$ \\
$-2.9$ & $-2$ & $-10^{19}$ & $-2^{63}$ \\
$9.3\times10^{18}$ & $2^{63}-1$ & $+\infty$ & $2^{63}-1$ \\
NaN & $0$ & & \\
\bottomrule
\end{sdtable}
\caption{stardust \code{astype(int64)}의 결과(복제본). $0$ 쪽으로 자르고, 범위 밖은 포화하고, NaN은 $0$이다.}
\label{tab:astype}
\end{table}

\code{bool}로 바꾸는 규칙은 \code{x != 0} 하나이다. 그래서 NaN은 참이고 $\pm0$은 거짓이다. \code{[nan, -0.0, 0.0]}은 \code{[True, False, False]}가 된다.

\subsection{Rounding to Integers}

실수를 가장 가까운 정수로 보내는 \code{round}의 동작은 언어마다 다르다. Python의 \code{round}와 numpy와 stardust의 \code{round}는 한가운데 값을 짝수로 보내서 \code{round(2.5)}가 $2$, \code{round(-2.5)}가 $-2$, \code{round(0.5)}가 $0$, \code{round(1.5)}가 $2$이다(stardust에서 확인). C의 \code{round}는 $0$에서 먼 쪽으로 보내서 \code{round(2.5)}가 $3$이고 C의 \code{rint}와 \code{nearbyint}는 현재 반올림 모드, 곧 보통은 짝수로의 반올림을 따른다. 같은 이름의 함수가 서로 다른 규칙을 따르므로 다른 언어의 코드를 옮길 때 이 차이가 조용히 결과를 바꾼다. binary64를 정수로 반올림하면 결과가 다시 binary64에 정확히 들어가므로 반올림 오차가 생기지 않는다. $|x|\ge2^{52}$이면 $x$ 자신이 답이다.

stardust의 \code{floor}, \code{ceil}, \code{trunc}, \code{round}는 C의 \code{std::floor}, \code{std::ceil}, \code{std::trunc}, \code{std::rint}이고 컴파일러는 이것을 각각 AArch64의 \code{FRINTM}, \code{FRINTP}, \code{FRINTZ}, \code{FRINTX} 명령 하나로 옮긴다. 결과는 부호를 지킨다. $-0.4$를 반올림하면 $-0$이고 $-0.5$를 $0$ 쪽으로 자르거나 올리면 역시 $-0$이다(\cref{tab:round-sd}). 정수 부분이 $0$인 음수의 결과가 $-0$인 것은 결과의 부호를 입력의 부호와 같게 두라는 표준의 규칙 때문이다.

\begin{table}[htbp]
\begin{sdtable}{colspec={lrrrrr}, row{1}={font=\small\bfseries}}
\toprule
$x$ & $2.5$ & $-2.5$ & $-0.4$ & $-0.5$ & $4503599627370495.5$ \\
\midrule
\code{round} & $2$ & $-2$ & $-0$ & $-0$ & $4503599627370496$ \\
\code{floor} & $2$ & $-3$ & $-1$ & $-1$ & $4503599627370495$ \\
\code{ceil} & $3$ & $-2$ & $-0$ & $-0$ & $4503599627370496$ \\
\code{trunc} & $2$ & $-2$ & $-0$ & $-0$ & $4503599627370495$ \\
\bottomrule
\end{sdtable}
\caption{stardust의 정수 반올림(복제본). 마지막 열은 $2^{52}-1/2$로 간격이 $1/2$인 마지막 이진 구간의 한가운데 값이다.}
\label{tab:round-sd}
\end{table}

명령이 없는 곳에서는 덧셈 하나로 짝수로의 반올림을 얻는 오래된 기법을 쓴다.

\begin{proposition}
\label{prop:magic-rint}
$C=1.5\cdot2^{52}$이라 하자. $|y|\le2^{51}$인 binary64 수 $y$에 대해 $(y\oplus C)\ominus C=\operatorname{rint}(y)$이다. 여기서 $\operatorname{rint}$는 한가운데를 짝수로 보내는 정수 반올림이다. 또 $\mathrm{bits}(y\oplus C)-\mathrm{bits}(C)=\operatorname{rint}(y)$이다.
\end{proposition}

\begin{proof}
$y+C\in[2^{52},2^{53}]$이다. 이진 구간 $[2^{52},2^{53})$의 간격은 $1$이므로 이 구간의 binary64 수는 정확히 정수들이고 유효숫자가 곧 그 정수이다. 따라서 $y+C$를 binary64로 반올림하는 것은 가장 가까운 정수로 보내는 것이고 한가운데에서 유효숫자가 짝수인 쪽은 정수가 짝수인 쪽이다. $C$가 짝수이므로 $\fl(y+C)=C+\operatorname{rint}(y)$이다($2^{53}$은 다음 구간의 첫 수이지만 정수이므로 같다). 그다음 뺄셈은 두 수가 모두 $[2^{52},2^{53}]$에 있으므로 \cref{lem:sterbenz}에 따라 정확하다. 마지막 주장은 \cref{prop:fp-bits-monotone}에서 이 구간의 연속한 정수가 연속한 비트에 대응한다는 데서 나온다.
\end{proof}

\begin{algorithm}[htbp]
\begin{algorithmic}[1]
\Require binary64 수 $y$
\Ensure $\operatorname{rint}(y)$, $0$의 부호를 지킨다
\State \textbf{if} $|y|\ge2^{52}$ 또는 $y$가 NaN \textbf{then} \Return $y$ \Comment{이미 정수, $\infty$, NaN}
\State $r\gets(y\oplus C)\ominus C$ \Comment{\cref{prop:magic-rint}, $C=1.5\cdot2^{52}$}
\State \Return \Call{CopySign}{$r$, $y$} \Comment{$-0.4$에서 $+0$ 대신 $-0$}
\end{algorithmic}
\caption{덧셈으로 하는 짝수로의 반올림}
\label{alg:rint-magic}
\end{algorithm}

이 기법은 정수 반올림보다 인자 축소에서 더 자주 쓰인다. stardust의 NEON $\exp$는 $x\cdot128/\log2$에 $C$를 더해 가장 가까운 정수 $N=128k+j$를 얻는다(\code{vmath.hpp}의 \code{kExpShift}). 덧셈 결과의 비트에는 $N$이 그대로 들어 있다. \cref{prop:magic-rint}의 마지막 주장에 따라 그 비트의 아래 $7$비트가 표의 칸 번호 $j$이다. 비트를 $45$비트 왼쪽으로 밀면 $C$의 비트는 $2^{64}$ 밖으로 밀려나고 $N\ll45=(k\ll52)+(j\ll45)$가 남는다. 표에는 $2^{j/128}$의 비트에서 $j\ll45$를 미리 뺀 값이 들어 있으므로 더하면 지수 필드에 $k$가 더해진 $2^{k+j/128}$이 된다(\cref{alg:split-exponent}의 \textproc{ScaleByPow2}). 덧셈 하나와 정수 연산 셋으로 반올림과 표 찾기와 $2^k$ 곱하기를 모두 끝낸다. 자세한 오차 해석은 13장의 몫이다. 소수 자리까지 반올림하는 \code{round(x, decimals)}는 사정이 다르다. 리터럴 \code{0.215}의 근삿값은 $0.21499999\cdots$이므로 올바른 답은 $0.21$이고 Python의 \code{round(0.215, 2)}는 정확한 십진 변환으로 이 답을 낸다. numpy의 \code{round(0.215, 2)}는 먼저 $100$을 곱하는데 그 곱이 반올림되어 정확히 $21.5$가 되고 짝수 쪽 $22$로 가서 $0.22$를 낸다(복제본의 numpy로 확인). 곱셈에 끼어든 반올림 한 번이 한가운데 값을 만든 것으로 \cref{sec:rounding}의 이중 반올림과 같은 모양이다. stardust의 \code{round}는 소수 자리 인자를 받지 않는다.

\section{Other Formats}
\label{sec:other-formats}

IEEE~754는 binary64 말고도 여러 형식을 정하고 하드웨어는 표준 밖의 형식도 쓴다(\cref{tab:formats}). 모두 \cref{def:fp-set}의 꼴이고 $p$와 $e_{\max}$만 다르다. 이 장의 정리들은 $p$, $e_{\min}$, $e_{\max}$와 거기서 나오는 상수($2^{-1074}$ 같은 것)를 바꾸어 그대로 성립한다. 형식 사이를 옮길 때 넓히는 쪽은 언제나 정확하다.

\begin{proposition}
\label{prop:fp-widening}
$p\le p'$, $e_{\min}'\le e_{\min}$, $e_{\max}\le e_{\max}'$이면 $\mathbb F(p,e_{\min},e_{\max})\subset\mathbb F(p',e_{\min}',e_{\max}')$이다.
\end{proposition}

\begin{proof}
$x=\pm m\cdot2^{e-p+1}$, $0\le m<2^p$, $e_{\min}\le e\le e_{\max}$라 하자. $m'=m\cdot2^{p'-p}<2^{p'}$로 두면 $x=\pm m'\cdot2^{e-p'+1}$이고 $e$는 $[e_{\min}',e_{\max}']$에 있다. 서브노멀 수도 $m$의 범위에 들어 있으므로 따로 다룰 것이 없다.
\end{proof}

\begin{table}[htbp]
\begin{sdtable}{colspec={lrrrl}, row{1}={font=\small\bfseries}}
\toprule
형식 & 비트 & $p$ & $e_{\max}$ & $u$ \\
\midrule
binary16 (half) & 16 & 11 & 15 & $2^{-11}\approx4.9\times10^{-4}$ \\
bfloat16 & 16 & 8 & 127 & $2^{-8}\approx3.9\times10^{-3}$ \\
binary32 (single) & 32 & 24 & 127 & $2^{-24}\approx6.0\times10^{-8}$ \\
binary64 (double) & 64 & 53 & 1023 & $2^{-53}\approx1.1\times10^{-16}$ \\
x87 확장 & 80 & 64 & 16383 & $2^{-64}\approx5.4\times10^{-20}$ \\
binary128 (quad) & 128 & 113 & 16383 & $2^{-113}\approx9.6\times10^{-35}$ \\
\bottomrule
\end{sdtable}
\caption{부동소수점 형식. $e_{\min}=1-e_{\max}$이다. bfloat16은 binary32의 위 $16$비트로 지수 범위가 같고 x87 확장 형식은 숨은 비트를 저장한다.}
\label{tab:formats}
\end{table}

짧은 형식은 기계 학습과 그래픽스에서 쓴다. binary16은 지수 범위가 좁아 가장 큰 유한 수가 $65504$이고 짝수로의 반올림에서는 $65520$ 이상이 넘친다. bfloat16은 binary32와 범위가 같은 대신 유효숫자가 $8$비트뿐이다. 신경망의 가중치처럼 오차를 견디는 자료라면 메모리와 대역폭을 반으로 줄이는 이득이 크다. 그러나 과학 계산의 대부분, 예컨대 조건수가 $10^6$인 선형계를 푸는 데는 binary32의 $u\approx6\times10^{-8}$조차 모자라다(11장의 섭동 이론에서 답의 상대오차가 대략 $\kappa u$이므로). 긴 형식은 정확도를 높이지만 binary128은 대부분의 하드웨어에서 소프트웨어로 흉내 내므로 binary64보다 수십 배 느리고 x87 확장 형식은 x86에만 있다. 넓은 형식을 거친 중간값이 \cref{thm:double-rounding}의 조건을 어기면 이중 반올림이 생긴다는 점도 보았다.

stardust의 dtype은 \code{bool}, \code{int64}, \code{float64}, \code{complex128} 넷뿐이고 실수는 binary64 하나이다. dtype 이름으로 \code{"float32"}를 주면 \code{TypeError}이다. numpy의 \code{float32}나 \code{float16} 배열은 받아들이되 가져오면서 binary64로 넓힌다. \cref{prop:fp-widening}에 따라 이 넓히기는 정확하다. \code{float32}의 \code{0.1}은 $0.10000000149011612$로, \code{float16}의 서브노멀 수 $2^{-23}$(\code{1e-7}을 반올림한 값)은 $1.1920928955078125\times10^{-7}$로 들어온다. \code{longdouble} 버퍼와 \code{uint64} 버퍼는 \code{TypeError}이다. 뒤의 것은 손실 없이 \code{int64}로 바꿀 수 없기 때문이다(복제본). 형식을 하나로 두면 모든 커널을 한 번만 짜고 한 번만 검증하면 되며(dtype이 늘면 커널과 정확도 계약이 그 수만큼 곱절로 는다) 승격 규칙이 한 줄로 끝난다(22장). Apple의 CPU에서 binary32를 쓰면 NEON 레지스터에 원소가 넷 들어가지만 그 이득이 과학 계산의 정확도를 포기할 만큼 크지 않다고 보았다. 더 높은 정밀도가 필요한 곳은 형식을 늘리지 않고 binary64 두 개로 만드는 double-double로 해결한다(12장).

\section{Compilers and Floating-Point Expressions}
\label{sec:compiler-fp}

\subsection{Associativity}

실수의 덧셈은 결합법칙을 만족하지만 부동소수점 덧셈은 그렇지 않다. $(2^{53}\oplus1)\oplus1$을 왼쪽부터 계산하면 첫 덧셈의 결과가 짝수로 반올림되어 $2^{53}$이 되고 다시 $1$을 더해도 $2^{53}$이다. 반면 $2^{53}\oplus(1\oplus1)$은 정확히 $2^{53}+2$이다. 그래서 컴파일러는 부동소수점 식의 계산 순서를 바꿀 수 없고 프로그래머가 적은 순서대로 연산을 해야 한다. 배열의 합처럼 순서가 결과를 바꾸는 연산에서 라이브러리는 순서를 스스로 정하고 그 순서를 약속해야 한다(29장의 쌍별 합).

많은 컴파일러에는 이 제약을 푸는 선택지가 있다. 이른바 \code{-ffast-math}는 결합법칙과 분배법칙이 성립하고 NaN과 무한대가 나오지 않는다고 가정하며 부호 있는 영을 구별하지 않는다. 그러면 합을 여러 갈래로 나누어 벡터화하는 등 더 빠른 코드를 만들 수 있다. 하지만 그 가정이 깨지는 입력에서는 결과를 예측할 수 없다. NaN이 나오지 않는다고 가정한 컴파일러는 \code{x != x} 검사를 거짓으로 바꾸어 없애 버릴 수 있다. 12장의 보상 합처럼 반올림 오차를 정확히 계산해 내는 알고리즘은 $(a+b)-a-b$처럼 대수적으로 $0$인 식을 지워 버리는 최적화 앞에서 무력해진다. stardust는 \code{-ffast-math}를 쓰지 않는다. 빌드 설정의 주석은 NaN과 무한대가 그 의미를 지켜야 한다는 점을 이유로 든다.

\subsection{Contraction and FMA}

이보다 미묘한 것이 \term{축약}{contraction}이다. 많은 프로세서에는 $a\cdot b+c$를 반올림 한 번으로 계산하는 fused multiply-add(FMA) 명령이 있다. FMA는 곱을 먼저 반올림하고 다시 더하는 것보다 정확하다. 그래서 C 표준은 컴파일러가 \code{a * b + c}를 FMA로 축약하는 것을 허락한다. 그러나 축약은 결과의 비트를 바꾼다. \cref{tab:fma}의 첫 예에서 $x=1+2^{-27}$이면 $x^2-1$의 참값은 $2^{-26}+2^{-54}$인데 곱을 먼저 반올림하면 $2^{-54}$가 사라지고 FMA는 참값을 그대로 낸다. 둘째 예는 거꾸로 FMA가 ``틀린'' 듯한 답을 낸다. \code{0.1*10}은 반올림으로 정확히 $1$이 되므로 $1$을 빼면 $0$이다. FMA는 \code{0.1}의 근삿값이 품은 오차를 그대로 드러내 $2^{-54}\approx5.55\times10^{-17}$을 낸다. 어느 쪽이든 축약 여부가 답을 바꾼다.

\begin{table}[htbp]
\begin{sdtable}{colspec={lll}, row{1}={font=\small\bfseries}}
\toprule
식 & 곱을 반올림한 뒤 더함 & FMA \\
\midrule
$x\cdot x-1$, $x=1+2^{-27}$ & $1.4901161193847656\times10^{-8}$ & $1.4901161249358807\times10^{-8}$ \\
$0.1\cdot10-1$ & $0$ & $5.551115123125783\times10^{-17}$ \\
\bottomrule
\end{sdtable}
\caption{축약이 바꾸는 결과(Python \code{math.fma}). 첫 줄의 FMA 값은 참값 $2^{-26}+2^{-54}$와 같고 둘째 줄의 FMA 값은 \code{0.1}의 근삿값과 $0.1$의 차이의 $10$배이다.}
\label{tab:fma}
\end{table}

축약이 문제 되는 까닭은 정확도에 있지 않고 재현성에 있다. 같은 식이 인라인된 위치나 벡터화 여부에 따라 축약되기도 하고 되지 않기도 하므로 같은 수식이 경로마다 다른 비트를 낼 수 있다. stardust는 \code{-ffp-contract=off}로 컴파일러의 축약을 끈다. 빌드 설정의 주석이 드는 예는 둘이다. 배열 $x$에 함수를 적용한 결과의 $i$번째 원소와 $x[i]$ 하나에 적용한 결과가 같아야 한다($f(x)[i]=f(x[i])$). 또 NEON 벡터 커널이 그 스칼라 정의와 같은 비트를 내야 한다. FMA가 필요한 곳에서는 코드가 \code{std::fma}나 \code{vfmaq\_f64}를 직접 적는다. 앞에서 본 \code{expm1}의 $\mathrm{fma}(s,t,s-1)$이 그런 예이다. 여기서 FMA는 컴파일러가 고른 것이 아니라 알고리즘의 일부이다. 같은 이유로 컴파일러가 같은 인자의 \code{sin}과 \code{cos}를 한 번의 \code{\_\_sincos\_stret} 호출로 합치는 것도 \code{-fno-builtin-sin}과 \code{-fno-builtin-cos}로 막고 합치고 싶은 곳에서는 코드가 직접 합친다.

stardust 자신에게도 FMA를 일부러 적은 곳이 있다. 복소수 곱 $(a_r+ia_i)(b_r+ib_i)$의 실수부를 $\mathrm{fma}(a_r,b_r,-a_ib_i)$로 계산하는 \code{cmul}(\code{scalar\_ops.hpp})이다. numpy의 arm64 빌드는 컴파일러의 기본 축약으로 \code{ar*br - ai*bi}를 이 꼴로 바꾼다. 그래서 numpy와 같은 비트를 내려면 같은 FMA를 손으로 적어야 한다. 두 꼴의 오차는 대칭이 아니다. 따로 계산하면 $(a_rb_r(1+\delta_1)-a_ib_i(1+\delta_2))(1+\delta_3)$이고 FMA는 $(a_rb_r-a_ib_i(1+\delta_2))(1+\delta_3)$이어서 첫 곱의 반올림이 사라진다. 무작위 복소수 $2\times10^4$쌍에서 두 꼴의 실수부는 $24.4\%$가 달랐고 numpy와 stardust는 FMA 꼴과 $100\%$ 같았다(복제본, \code{figures/10/review\_sd.py}). 역수 \code{reciprocal}의 분모도 같은 이유로 FMA를 적는다. 코드 주석에 따르면 FMA 없이는 $9\%$의 결과가 numpy와 달랐다.

\subsection{The Bit Policy}

이런 선택은 stardust가 결과의 비트에 대해 무엇을 약속하는지와 이어진다. 약속은 이렇다. 같은 기계에서 같은 스레드 수로 실행하면 실행할 때마다 같은 비트가 나온다. numpy나 scipy와 비트까지 같다는 약속은 하지 않는다. 대신 반올림 수준의 차이는 Arb로 구한 참값에 대한 정확도 계약(``$k$ ULP 이내'') 안에 있으면 된다. 반올림이 끼지 않는 동작, 곧 dtype 규칙과 정수 연산, 정렬 순서, NaN과 $\pm0$의 규칙은 numpy와 같게 둔다. 이 약속을 지키려면 이 절의 세 가지가 필요하다. 계산 순서를 코드가 정하고(결합법칙), 축약을 코드가 정하고(FMA), 특수값의 의미를 지킨다(fast-math 없음). 더 빠르거나 더 정확한 계산 순서를 고를 자유는 남는다. 다만 그 순서는 입력의 길이와 배치에만 달려 있어야 하고 실행 시점의 사정(스레드의 도착 순서 같은 것)에 달려 있으면 안 된다.
\section{Testing Floating-Point Code}
\label{sec:fp-testing}

부동소수점 커널의 시험은 무작위 입력만으로는 모자라다. $\pm0$, $\pm\infty$, NaN, 서브노멀 수, 정규 수의 경계 $2^{-1022}$, 가장 큰 유한 수, $2$의 거듭제곱, 그리고 정수 변환의 경계 $\pm2^{63}$과 $2^{53}$처럼 이 장의 규칙이 달라지는 지점은 무작위로 뽑으면 거의 나오지 않는다. 그래서 stardust의 시험은 여러 갈래로 짠다(\cref{tab:fp-tests}). 첫째는 특수값의 모든 쌍이다. 이항 연산이면 특수값 목록의 곱집합 전체를 넣고 결과를 비트로 견준다. 이때 \code{==}로 견주면 $-0$과 $+0$을 구별하지 못하고 NaN끼리는 같지 않다고 나오므로 바이트를 견주거나(\code{tobytes}) 부호 비트를 따로 견준다. 둘째는 경계 근처의 촘촘한 표본이다. 결과가 $0$이 되는 점($\log$의 $1$, $\cos$의 $\pi/2$의 홀수배)과 공식이 바뀌는 문턱 근처를 촘촘히 뽑고 \code{nextafter}로 경계의 양옆 몇 ULP를 걷는다. 셋째는 넓은 범위의 무작위 입력으로 $e^{U}$처럼 지수를 고르게 뽑아야 이진 구간마다 비슷한 수의 표본이 들어간다.

특수값 목록에는 조용한 NaN 하나만이 아니라 신호 NaN, 페이로드가 있는 NaN, 부호 비트가 켜진 NaN도 넣어야 한다. stardust의 특수값 목록에는 이것들이 빠져 있어서 \cref{sec:special-values}에서 본 신호 NaN의 경로 차이가 시험에 걸리지 않았다. 비교 자체는 \cref{alg:same-bits}처럼 짠다.

\begin{algorithm}[htbp]
\begin{algorithmic}[1]
\Require 이항 연산 $f$와 그 참 정의 $g$, 특수값 목록 $L$ (비트로 적는다)
\Function{SameBits}{$x$, $y$} \Comment{$\pm0$과 NaN의 비트까지}
  \State \Return $\mathrm{bits}(x)=\mathrm{bits}(y)$
\EndFunction
\ForAll{$a\in L$, $b\in L$}
  \For{$n\in\{1,2,3,5,17\}$} \Comment{스칼라 경로, 벡터 경로, 꼬리}
    \State $A\gets$ 길이 $n$의 배열, 모든 원소 $a$; $B\gets$ 같은 꼴로 $b$
    \State \textbf{assert} 모든 $i$에서 \Call{SameBits}{$f(A,B)_i$, $g(a,b)$}
  \EndFor
\EndFor
\end{algorithmic}
\caption{특수값 곱집합 시험. 길이를 바꾸어 벡터 경로와 스칼라 꼬리를 모두 지나게 한다.}
\label{alg:same-bits}
\end{algorithm}

견줄 상대도 정해야 한다. 기본 연산과 \code{nextafter}, \code{maximum}, 정수 변환처럼 결과가 하나로 정해진 연산은 참값과 비트까지 같아야 한다. 초월 함수처럼 올바른 반올림을 약속하지 않는 연산은 Arb의 구간 산술로 구한 참값과의 ULP 거리(\cref{alg:ulp-distance})를 재고 계약(``$k$ ULP 이내'')과 견준다. ULP 거리를 재는 도구 자신도 시험이 필요하다. stardust의 시험은 $\pm0$ 사이가 $0$, 이웃한 수가 $1$, $\pm2^{-1074}$ 사이가 $2$, NaN 둘이 $0$, NaN과 수가 표지값, $2$의 거듭제곱을 넘는 한 걸음이 $1$, $-\infty$와 $+\infty$ 사이가 $2^{64}-2^{53}$임을 확인한다.

\begin{table}[htbp]
\begin{sdtable}{colspec={X[1.4,l]X[3,l]X[1.6,l]}, row{1}={font=\small\bfseries}, cells={font=\small}}
\toprule
층 & 입력과 기준 & stardust의 예 \\
\midrule
특수값의 곱집합 & $\pm0$, $\pm\infty$, NaN, $\pm2^{-1074}$, $10^{-310}$, $\pm1$, $1/2$, $2$, $10^{308}$의 모든 쌍, 바이트 비교 & \code{pow}의 특수값, \code{maximum}과 \code{minimum}의 $\pm0$과 NaN \\
경계 근처 & 함수가 갈리는 점 근처의 촘촘한 표본, 경계 양옆 몇 ULP & NEON \code{exp}, \code{log}, \code{cos}($0$, $1$, $\pi/2$의 홀수배 근처) \\
넓은 무작위 & 지수를 고르게 뽑은 입력, Arb 참값과의 ULP 거리 & 원소 함수의 정확도 계약 \\
변환의 경계 & $\pm2^{63}$, $2^{53}$, NaN, $\pm\infty$, 소수 부분 & \code{astype(int64)}의 포화 \\
도구 자신 & ULP 거리의 알려진 값 & \code{ulp\_distance} \\
\bottomrule
\end{sdtable}
\caption{부동소수점 코드의 시험 층. 예는 stardust의 \code{test\_ufunc.py}, \code{test\_astype\_}\allowbreak\code{assignment.py}, \code{test\_accuracy\_tools.py}에서 가져왔다.}
\label{tab:fp-tests}
\end{table}

이 장에서 직접 확인한 것은 \cref{tab:binary64}의 상수(정확한 유리수), \cref{fig:relerr}의 반올림 오차 한계, \cref{fig:bias}의 치우침, 서막의 $0.1+0.2$, \cref{tab:special-sd}와 \cref{tab:astype}의 특수값과 변환 결과, 서브노멀 곱셈의 시간, \cref{tab:fma}의 축약 예이다. 수치는 모두 Python의 \code{fractions}와 \code{decimal}로 정확히 계산했고 stardust의 결과는 따로 빌드한 복제본에서 얻었다.

\section{APIs in This Chapter}
\label{sec:fp-apis}

\cref{tab:fp-apis}에 이 장의 규칙을 드러내는 stardust API를 모았다. 배열에 싣는 방법과 인자 규칙은 3권이 다룬다.

\begin{table}[htbp]
\begin{sdtable}{colspec={X[2,l]X[3,l]X[1.5,l]}, row{1}={font=\small\bfseries}}
\toprule
API & 이 장에서 다룬 것 & 설명한 곳 \\
\midrule
\code{finfo(float64)} & \code{eps}, \code{max}, \code{min}, \code{smallest\_normal} & \cref{tab:binary64}, 22장 \\
\code{nextafter} & 비트 패턴의 걸음 & \cref{alg:nextafter}, 27장 \\
\code{copysign}, \code{signbit} & 부호 있는 영의 부호 읽기 & \cref{sec:special-values}, 27장 \\
\code{maximum}, \code{minimum}, \code{max}, \code{min} & IEEE \code{maximum}, $-0<+0$, NaN 전파 & \cref{alg:fmax-nan}, 29장 \\
\code{nanmax}, \code{nanmin} & IEEE \code{maximumNumber} & \cref{alg:fmax-nan}, 29장 \\
\code{isnan}, \code{isinf}, \code{isfinite} & 비트로 하는 판정 & \cref{alg:bit-predicates} \\
\code{astype(int64)} & $0$ 쪽 자름, 포화, NaN $\to0$ & \cref{alg:float-to-int}, 22장 \\
\code{round}, \code{trunc}, \code{floor}, \code{ceil} & 짝수로의 반올림과 세 방향, $-0$ & \cref{tab:round-sd}, \cref{alg:rint-magic}, 27장 \\
\code{astype(float64)}, \code{asarray} & 정수의 올바른 반올림, 좁은 형식 넓히기 & \cref{alg:int-to-float}, \cref{prop:fp-widening} \\
\code{mean}(정수 배열) & 정확한 정수합 뒤 반올림 & \cref{sec:conversion}, 29장 \\
\bottomrule
\end{sdtable}
\caption{이 장의 규칙을 쓰는 stardust API. stardust에는 numpy의 \code{fmax}, \code{spacing}, \code{frexp}, \code{ldexp}, \code{isclose}가 없다.}
\label{tab:fp-apis}
\end{table}

\section{Concluding Remarks}

부동소수점 수는 실수를 흉내 내는 불완전한 장치가 아니라 정확한 규칙을 따르는 수 체계이다. 수들의 집합과 간격이 정해져 있다. 기본 연산은 정확한 결과를 그 집합으로 반올림한 것이고 그 반올림의 오차는 \cref{thm:standard-model}로 묶인다. 특수값도 각자의 규칙을 따르므로 계산의 결과는 원칙적으로 언제나 예측할 수 있다. 구현에서 이 규칙은 비트 수준의 도구가 되었다. 비트 패턴이 단조이므로 이웃한 수를 찾는 일과 두 수의 거리를 재는 일이 정수 덧셈과 뺄셈이 되었다. 표준의 \code{maximum}과 \code{FCVTZS}가 하드웨어 명령과 맞아떨어지니 규칙을 지키면서도 벡터화할 수 있었다. 거꾸로 컴파일러에게 맡기면 결합법칙과 축약이 비트를 바꾸므로 계산 순서와 FMA의 위치는 코드가 스스로 정해야 했다. 다음 장에서는 이 규칙을 출발점으로 삼아 연산이 여러 번 이어질 때 오차가 어떻게 쌓이고 퍼지는지를 해석한다. 12장에서는 반올림 오차 자체를 정확히 계산해 내는 방법으로 정확도를 되찾는다.

부동소수점 산술 전반을 알고 싶다면 Goldberg의 개관\cite{goldberg1991}이 고전적인 입문이다. Overton\cite{overton2001}은 짧고 친절한 교과서이며 Muller 등\cite{muller2018handbook}은 이 분야의 표준적인 참고서이다. Sterbenz의 보조정리는 \cite{sterbenz1974}에서 처음 정리되었고 이중 반올림의 조건은 Figueroa\cite{figueroa1995}가 밝혔다. 표준 자체는 \cite{ieee754-2019}이다.

\chapterreferences
